\documentclass[12pt,a4paper]{ctexart}

% 页面设置
\usepackage[margin=2.5cm]{geometry}

% 数学相关
\usepackage{amsmath,amssymb,amsfonts}
\usepackage{mathtools}

% 表格相关
\usepackage{booktabs}
\usepackage{longtable}
\usepackage{array}

% 图形和颜色
\usepackage{graphicx}
\usepackage{xcolor}

% 列表
\usepackage{enumitem}

% 超链接
\usepackage{hyperref}
\hypersetup{
    colorlinks=true,
    linkcolor=blue,
    urlcolor=blue
}

% 页眉页脚
\usepackage{fancyhdr}
\pagestyle{fancy}
\fancyhf{}
\fancyhead[L]{优化算法：从理论到大规模实践 \quad 课程习题集}
\fancyhead[R]{\thepage}
\renewcommand{\headrulewidth}{0.4pt}

% 自定义环境
\usepackage{tcolorbox}
\tcbuselibrary{skins,breakable}

% 习题环境
\newcounter{problem}[section]
\renewcommand{\theproblem}{\arabic{problem}}

\newenvironment{problem}[1][]{%
    \refstepcounter{problem}%
    \par\medskip\noindent%
    \textbf{题 \theproblem}%
    \ifx&#1&\else~\textbf{（#1）}\fi%
    \quad%
}{\par\medskip}

\newenvironment{subproblems}{%
    \begin{enumerate}[label=(\alph*),leftmargin=2em]%
}{%
    \end{enumerate}%
}

% 提示环境
\newtcolorbox{hint}[1][提示]{
    colback=yellow!5,
    colframe=orange!50!black,
    title=#1,
    fonttitle=\bfseries\small,
    breakable,
    boxrule=0.5pt
}

\title{\Huge\bfseries 优化算法：从理论到大规模实践\\[0.5em]
\Large 课程习题集}
\author{}
\date{}

\begin{document}

\maketitle
\thispagestyle{empty}

\vspace{1cm}
\begin{center}
\textbf{说明}：本习题集按课程15讲的内容分章节编排，每章包含概念题、计算题、证明题和思考题。\\
题目难度：$\star$ 基础 \quad $\star\star$ 中等 \quad $\star\star\star$ 挑战
\end{center}

\newpage
\tableofcontents
\newpage

%========================================
% 第1章
%========================================
\section{第1讲：优化的边界（计算复杂性）}

\begin{problem}[概念题 $\star$]
请简述P类问题和NP类问题的定义，并解释"P $\subseteq$ NP"为什么成立。
\end{problem}

\begin{problem}[概念题 $\star$]
解释以下概念之间的区别：
\begin{subproblems}
    \item 可计算问题与不可计算问题
    \item P类问题与NP类问题
    \item NP-Hard问题与NP-Complete问题
\end{subproblems}
\end{problem}

\begin{problem}[概念题 $\star\star$]
用停机问题（Halting Problem）的不可判定性来解释"优化的上限"。请给出停机问题不可判定的简要证明思路（对角化论证）。
\end{problem}

\begin{problem}[计算题 $\star$]
考虑旅行商问题（TSP），$n$个城市。
\begin{subproblems}
    \item 暴力枚举算法的时间复杂度是多少？
    \item 当$n=15$时，暴力枚举大约需要计算多少条不同的路径？
    \item 使用动态规划（Held-Karp算法）可以将时间复杂度降到多少？
\end{subproblems}
\end{problem}

\begin{problem}[概念题 $\star\star$]
什么是多项式时间归约（polynomial-time reduction）？请说明如何将3-SAT问题归约到独立集问题（Independent Set），并解释这如何证明独立集问题是NP-Complete的。
\end{problem}

\begin{problem}[思考题 $\star\star$]
课程提出"上限-下限"框架：上限是不可计算，下限是暴力枚举$O(2^n)$。请解释：
\begin{subproblems}
    \item 为什么说"科研的本质是发现结构、发展针对性方法"？
    \item 请用背包问题为例，说明如何通过发现不同层次的结构（如贪心近似、DP精确解、FPTAS）在上限和下限之间找到好的算法。
\end{subproblems}
\end{problem}

\begin{problem}[思考题 $\star\star\star$]
GPU并行可以改变我们对复杂度的看法。假设有一个$O(n^2)$的算法A可以完全并行化（每步独立），和一个$O(n\log n)$的算法B本质上是顺序的。
\begin{subproblems}
    \item 在拥有$p$个并行处理器的情况下，算法A的并行时间复杂度是多少？
    \item 在什么条件下（$n$和$p$的关系），并行化的A比顺序的B更快？
    \item 这对"最优算法"的概念有何启示？
\end{subproblems}
\end{problem}

\begin{problem}[计算题 $\star\star$ \textbf{归约手算：图着色$\to$3-SAT}]
考虑图$G$有3个节点$\{v_1, v_2, v_3\}$和3条边$\{(v_1,v_2), (v_2,v_3), (v_1,v_3)\}$（完全图$K_3$）。使用3种颜色$\{R,G,B\}$的图着色问题。

将3-着色问题编码为SAT：为每个节点$v_i$和颜色$c$引入布尔变量$x_{i,c}$（表示"$v_i$被染为颜色$c$"）。
\begin{subproblems}
    \item \textbf{至少一种颜色}：为每个节点写一个子句，保证它至少被染一种颜色。例如$v_1$对应$(x_{1,R} \lor x_{1,G} \lor x_{1,B})$。写出全部3个子句。
    \item \textbf{至多一种颜色}：对每个节点，任意两种颜色不能同时为真。例如$v_1$不能同时为$R$和$G$：$(\neg x_{1,R} \lor \neg x_{1,G})$。写出全部$3 \times \binom{3}{2} = 9$个子句。
    \item \textbf{相邻不同色}：对每条边和每种颜色，两个端点不能同色。例如边$(v_1,v_2)$颜色$R$：$(\neg x_{1,R} \lor \neg x_{2,R})$。写出全部$3 \times 3 = 9$个子句。
    \item 共多少个变量？多少个子句？手动求解该SAT实例（找到所有满足赋值），验证它们恰好对应$K_3$的所有合法3-着色方案。$K_3$有多少种合法着色？
    \item 如果改为$K_4$（4个节点的完全图，仍用3种颜色），变量数和子句数分别变为多少？该SAT实例是否可满足？为什么？
\end{subproblems}
\end{problem}

\begin{problem}[计算题 $\star\star$ \textbf{归约手算：子集和$\to$SAT}]
子集和问题：给定集合$S = \{3, 5, 7\}$和目标$T = 8$，是否存在子集的和恰好为$T$？

用布尔变量$x_i$（$x_i = 1$表示选第$i$个元素）。
\begin{subproblems}
    \item 写出约束：$3x_1 + 5x_2 + 7x_3 = 8$。枚举所有$2^3 = 8$种赋值，哪些满足？
    \item 将整数约束转化为二进制表示的逻辑电路（加法器），再将电路转化为CNF子句（Tseitin变换）。画出$3x_1 + 5x_2$的二进制加法电路（$3 = 011_2$, $5 = 101_2$, 和最多$1000_2$）。
    \item 对一个全加器（3输入1位加法：$a + b + c_{\text{in}} = 2c_{\text{out}} + s$），写出$s$和$c_{\text{out}}$的逻辑表达式，并转化为CNF（Tseitin变换：引入中间变量，每个门$\leq 3$个子句）。
    \item 这个归约说明了什么？（子集和是NP-Complete，因此它可以被编码为SAT。）
\end{subproblems}
\end{problem}

\begin{problem}[计算题 $\star\star\star$ \textbf{Held-Karp DP手算}]
5个城市的TSP，距离矩阵：
\[
D = \begin{pmatrix}
0 & 3 & 4 & 2 & 7 \\
3 & 0 & 4 & 6 & 3 \\
4 & 4 & 0 & 5 & 8 \\
2 & 6 & 5 & 0 & 6 \\
7 & 3 & 8 & 6 & 0
\end{pmatrix}
\]
从城市0出发。定义$dp[S][j]$为从城市0出发，恰好经过集合$S$中所有城市，最后停在城市$j$的最短路径长度。
\begin{subproblems}
    \item 初始化：$dp[\{j\}][j] = D[0][j]$，$j = 1,2,3,4$。写出这4个值。
    \item 填写$|S|=2$的所有$dp$值：$dp[\{j,k\}][k] = \min_j(dp[\{j\}][j] + D[j][k])$。共有$\binom{4}{2} \times 2 = 12$个值。
    \item 填写$|S|=3$的$dp$值。
    \item 填写$|S|=4$的$dp$值。
    \item 最终答案：$\min_j(dp[\{1,2,3,4\}][j] + D[j][0])$。最优路径是什么？
    \item 暴力枚举需要检查$(5-1)!/2 = 12$条路径。Held-Karp需要填写多少个$dp$值？计算复杂度$O(2^n n^2)$在$n=5$时是多少？
\end{subproblems}
\end{problem}

\begin{problem}[计算题 $\star\star$ \textbf{贪心近似比手算}]
考虑集合覆盖问题。全集$U = \{1,2,3,4,5,6\}$，子集族：
\[
S_1 = \{1,2,3\}, \quad S_2 = \{2,4\}, \quad S_3 = \{3,5\}, \quad S_4 = \{4,5,6\}, \quad S_5 = \{1,6\}
\]
每个子集成本为1。
\begin{subproblems}
    \item 贪心算法：每步选覆盖最多未覆盖元素的集合。写出每一步的选择过程和覆盖进展。贪心解用了几个集合？
    \item 最优解是什么？用了几个集合？
    \item 计算近似比$\text{ALG}/\text{OPT}$。贪心集合覆盖的理论近似比上界是$H(|U|) = \sum_{i=1}^6 1/i$。计算此值。
    \item 将集合覆盖写成整数规划和对应的LP松弛。LP松弛的最优值是多少？（提示：可以赋分数解。）
\end{subproblems}
\end{problem}

%========================================
% 第2章
%========================================

% --- 按章节归并的补充题 ---
\subsection*{第1讲：计算复杂性}

\begin{problem}[$\star$]
\textbf{暴力 vs 贪心 vs 2-opt：TSP实验}

生成$n$个二维平面上的随机点作为城市。分别实现：
\begin{subproblems}
    \item \textbf{暴力搜索}：遍历所有排列，找最短路径。用 \texttt{itertools.permutations} 即可。
    \item \textbf{贪心}：每次走到最近的未访问城市。
    \item \textbf{2-opt}：从贪心解出发，反复尝试翻转子路径，直到无法改进。
\end{subproblems}
分别在$n=8,10,12$上运行，记录运行时间和路径长度。画表格比较三者。$n=12$时暴力法要多久？

\begin{hint}
2-opt的核心循环：对所有$(i,j)$对，检查将路径$[i..j]$翻转后总长度是否减小。
\end{hint}
\end{problem}

\begin{problem}[$\star$]
\textbf{SAT求解器体验}

用Python实现一个暴力SAT求解器：输入CNF公式（子句列表），枚举所有$2^n$种变量赋值，检查是否满足。
\begin{subproblems}
    \item 对以下CNF公式求解：$(x_1 \lor x_2) \land (\neg x_1 \lor x_3) \land (\neg x_2 \lor \neg x_3)$。
    \item 随机生成含$n$个变量、$m$个子句（每个子句3个文字）的3-SAT实例。测试$n=10,15,20$时的求解时间。
    \item 固定$n=20$，画出"可满足概率"随$m/n$比值变化的曲线（子句密度相变现象）。
\end{subproblems}
\end{problem}

%--- 第2讲 ---
\section{第2讲：凸优化与KKT}

\begin{problem}[概念题 $\star$]
给出凸集和凸函数的定义。证明：如果$f$是凸函数，则$f$的所有下水平集$\{x : f(x) \leq \alpha\}$是凸集。
\end{problem}

\begin{problem}[计算题 $\star$]
判断以下函数是否为凸函数，并说明理由（可使用二阶条件）：
\begin{subproblems}
    \item $f(x) = x^4 - 2x^2 + 1$，$x \in \mathbb{R}$
    \item $f(x_1, x_2) = x_1^2 + x_2^2 + x_1 x_2$
    \item $f(x) = \log(1 + e^x)$，$x \in \mathbb{R}$
    \item $f(X) = -\log\det(X)$，$X \in \mathbb{S}_{++}^n$（正定矩阵集）
\end{subproblems}
\end{problem}

\begin{problem}[概念题 $\star\star$]
解释Jensen不等式，并用它证明：对于凸函数$f$和随机变量$X$，有$f(\mathbb{E}[X]) \leq \mathbb{E}[f(X)]$。给出一个在机器学习中应用Jensen不等式的例子。
\end{problem}

\begin{problem}[证明题 $\star\star$]
证明凸优化的基本定理：如果$f$是凸函数，$X$是凸集，则$f$在$X$上的任何局部最小值都是全局最小值。
\end{problem}

\begin{problem}[计算题 $\star\star$]
考虑以下凸优化问题：
\[
\min_{x \in \mathbb{R}^2} \quad x_1^2 + 2x_2^2 \quad \text{s.t.} \quad x_1 + x_2 = 1, \quad x_1 \geq 0, \quad x_2 \geq 0
\]
\begin{subproblems}
    \item 写出该问题的Lagrangian函数$L(x, \lambda, \nu)$。
    \item 写出KKT条件。
    \item 求解最优解$x^*$和对应的对偶变量$\lambda^*, \nu^*$。
    \item 验证互补松弛条件$\lambda_i g_i(x^*) = 0$是否成立。
\end{subproblems}
\end{problem}

\begin{problem}[计算题 $\star\star$]
考虑Lagrangian对偶问题。对于
\[
\min_{x} \quad \frac{1}{2}x^\top Q x + c^\top x \quad \text{s.t.} \quad Ax = b
\]
其中$Q \succ 0$（正定）。
\begin{subproblems}
    \item 写出Lagrangian函数。
    \item 对$x$求极小，得到关于对偶变量$\nu$的对偶函数$g(\nu)$。
    \item 写出对偶问题，并验证强对偶性成立。
\end{subproblems}
\end{problem}

\begin{problem}[概念题 $\star\star$]
解释对偶变量$\lambda^*$的经济学含义（"影子价格"）。具体来说，如果约束从$g_i(x) \leq 0$变为$g_i(x) \leq \epsilon$，最优值$p^*$如何变化？
\end{problem}

\begin{problem}[证明题 $\star\star\star$]
\textbf{弱对偶定理}：证明对于一般的优化问题，对偶最优值$d^*$满足$d^* \leq p^*$（原问题最优值）。
\begin{hint}
从Lagrangian的定义出发，利用$\lambda \geq 0$和约束$g(x) \leq 0$。
\end{hint}
\end{problem}

%========================================
% 第3章
%========================================

% --- 按章节归并的补充题 ---
\subsection*{一、凸性判定与性质（第2讲）}

\begin{problem}[$\star$]
判断以下集合是否为凸集，如果不是请给出反例：
\begin{subproblems}
    \item $\{x \in \mathbb{R}^2 : x_1^2 + x_2^2 \leq 1\}$（单位圆盘）
    \item $\{x \in \mathbb{R}^2 : x_1^2 + x_2^2 = 1\}$（单位圆周）
    \item $\{x \in \mathbb{R}^2 : x_1 \geq 0, \ x_1 x_2 \geq 1\}$
    \item $\{X \in \mathbb{R}^{n\times n} : X = X^\top, \ X \succeq 0\}$（半正定锥）
\end{subproblems}
\end{problem}

\begin{problem}[$\star\star$]
设$f:\mathbb{R}^n \to \mathbb{R}$是凸函数。证明其上境图（epigraph）$\text{epi}(f) = \{(x,t) : f(x) \leq t\}$是凸集。反之，如果$\text{epi}(f)$是凸集，则$f$是凸函数。
\end{problem}

\begin{problem}[$\star\star$]
证明以下凸性保持运算：
\begin{subproblems}
    \item 如果$f_1, f_2$都是凸函数，则$g(x) = \max\{f_1(x), f_2(x)\}$也是凸函数。
    \item 如果$f(x,y)$关于$(x,y)$联合凸，则$g(x) = \inf_y f(x,y)$关于$x$是凸函数（假设下确界可以取到）。
    \item 如果$f$是凸函数，$A$是矩阵，$b$是向量，则$g(x) = f(Ax+b)$也是凸函数。
\end{subproblems}
\end{problem}

\begin{problem}[$\star$]
计算以下函数的Hessian矩阵，并判定凸性：
\begin{subproblems}
    \item $f(x,y) = e^{x+3y-0.1} + e^{x-3y-0.1} + e^{-x-0.1}$
    \item $f(x,y) = x^2/y$，定义域$y > 0$
    \item $f(x) = \|Ax - b\|_2^2 + \lambda \|x\|_2^2$，其中$A \in \mathbb{R}^{m \times n}$，$\lambda > 0$
\end{subproblems}
\end{problem}

%--- KKT条件手算（对应第2讲）---
\subsection*{二、KKT条件手算（第2讲）}

\begin{problem}[$\star\star$]
\textbf{(Boyd 5.1改编)} 求解以下问题的KKT条件并找到最优解：
\[
\min_{x \in \mathbb{R}^2} \quad (x_1 - 1)^2 + (x_2 - 2)^2 \quad \text{s.t.} \quad x_1 + x_2 \leq 2, \quad x_1, x_2 \geq 0
\]
要求写出完整的KKT条件（稳定性、原始可行性、对偶可行性、互补松弛），并求解$x^*$和$\lambda^*$。
\end{problem}

\begin{problem}[$\star\star$]
\textbf{(水池注水问题/Water-Filling)} 求解：
\[
\max_{x \in \mathbb{R}^n} \quad \sum_{i=1}^n \log(\alpha_i + x_i) \quad \text{s.t.} \quad \sum_{i=1}^n x_i = 1, \quad x_i \geq 0
\]
其中$\alpha_i > 0$为给定常数。
\begin{subproblems}
    \item 写出KKT条件。
    \item 证明最优解形如$x_i^* = \max(\nu - \alpha_i, 0)$，其中$\nu$由$\sum_i x_i^* = 1$确定。
    \item 当$n=3$，$\alpha = (1, 2, 4)$时，具体计算$x^*$和$\nu$。
\end{subproblems}
\end{problem}

\begin{problem}[$\star\star$]
\textbf{(SVM对偶)} 考虑线性SVM的原问题：
\[
\min_{w,b} \quad \frac{1}{2}\|w\|^2 \quad \text{s.t.} \quad y_i(w^\top x_i + b) \geq 1, \quad i=1,\ldots,m
\]
\begin{subproblems}
    \item 写出Lagrangian $L(w,b,\alpha)$。
    \item 对$w$和$b$分别求导令其为零，得到$w$和$b$关于$\alpha$的表达式。
    \item 代入得到对偶问题。
    \item 利用互补松弛$\alpha_i[y_i(w^\top x_i + b) - 1] = 0$解释"支持向量"。
    \item 给定数据点$x_1 = (1,0)$, $y_1=+1$; $x_2=(0,1)$, $y_2=+1$; $x_3=(-1,0)$, $y_3=-1$。手算最优$w^*$, $b^*$和$\alpha^*$。
\end{subproblems}
\end{problem}

\begin{problem}[$\star\star\star$]
\textbf{(Boyd 5.6改编)} 考虑带二次约束的二次规划（QCQP）：
\[
\min_{x \in \mathbb{R}^n} \quad \frac{1}{2}x^\top P_0 x + q_0^\top x \quad \text{s.t.} \quad \frac{1}{2}x^\top P_i x + q_i^\top x + r_i \leq 0, \quad i=1,\ldots,m
\]
其中所有$P_i \succeq 0$。
\begin{subproblems}
    \item 写出Lagrangian和对偶函数$g(\lambda)$。
    \item 证明对偶函数$g(\lambda)$是凹函数。
    \item 在一维情形$n=1$下，取$P_0=2$, $q_0=-2$, $P_1=2$, $q_1=0$, $r_1=-1$（即$\min (x-1)^2$ s.t. $x^2 \leq 1$），手算原问题和对偶问题的最优值，验证强对偶。
\end{subproblems}
\end{problem}

\begin{problem}[$\star\star$]
\textbf{(投影问题)} 求点$a = (3,4)$到以下凸集的投影（即集合中离$a$最近的点）：
\begin{subproblems}
    \item 单位球$\{x : \|x\|_2 \leq 1\}$
    \item 半空间$\{x : 2x_1 + x_2 \leq 3\}$
    \item 仿射集$\{x : x_1 + x_2 = 1\}$
    \item 单纯形$\{x : x_1 + x_2 = 1, x_1 \geq 0, x_2 \geq 0\}$
\end{subproblems}
用KKT条件求解每个子问题。
\end{problem}

%--- LP对偶手算（对应第3讲）---
\subsection*{第2讲：凸优化与KKT}

\begin{problem}[$\star\star$]
\textbf{用CVXPY求解并验证KKT}

考虑投资组合优化问题：
\[
\min_w \quad w^\top \Sigma w \quad \text{s.t.} \quad \mu^\top w \geq r_{\min}, \quad \mathbf{1}^\top w = 1, \quad w \geq 0
\]
其中$\Sigma$是$5\times 5$的协方差矩阵，$\mu$是期望收益向量。
\begin{subproblems}
    \item 随机生成$\Sigma$（半正定）和$\mu$，用CVXPY求解。
    \item 从求解结果中提取对偶变量，手动验证KKT条件（梯度条件、互补松弛）。
    \item 改变$r_{\min}$的值，画出有效前沿（efficient frontier）：横轴标准差，纵轴期望收益。
    \item 验证"影子价格"：$r_{\min}$增加$\epsilon$后，最优目标值变化是否约为$\lambda^* \cdot \epsilon$？
\end{subproblems}
\end{problem}

%--- 第3讲 ---
\section{第3讲：LP对偶与近似算法}

\begin{problem}[概念题 $\star$]
写出线性规划标准形及其对偶问题：
\begin{subproblems}
    \item 原问题：$\min c^\top x$ s.t. $Ax = b$, $x \geq 0$
    \item 对偶问题是什么？
    \item 写出LP对偶的机械化构造规则。
\end{subproblems}
\end{problem}

\begin{problem}[计算题 $\star\star$]
考虑以下线性规划问题：
\[
\max \quad 3x_1 + 5x_2 \quad \text{s.t.} \quad x_1 \leq 4, \quad 2x_2 \leq 12, \quad 3x_1 + 5x_2 \leq 25, \quad x_1, x_2 \geq 0
\]
\begin{subproblems}
    \item 写出该问题的对偶问题。
    \item 利用互补松弛条件和对偶可行性，求解原问题和对偶问题的最优解。
\end{subproblems}
\end{problem}

\begin{problem}[证明题 $\star\star$]
证明LP的强对偶定理（假设原问题有有界最优解）：原问题最优值等于对偶问题最优值，即$p^* = d^*$。
\begin{hint}
可以使用超平面分离定理或Farkas引理。
\end{hint}
\end{problem}

\begin{problem}[计算题 $\star\star$]
\textbf{LP松弛与近似}。考虑最小顶点覆盖问题（Minimum Vertex Cover）：给定图$G=(V,E)$，其整数规划形式为：
\[
\min \sum_{v \in V} x_v \quad \text{s.t.} \quad x_u + x_v \geq 1, \forall (u,v)\in E, \quad x_v \in \{0,1\}
\]
\begin{subproblems}
    \item 写出其LP松弛（将$x_v \in \{0,1\}$放宽为$0 \leq x_v \leq 1$）。
    \item 描述LP松弛舍入算法：$x_v^{LP} \geq 1/2$则取$x_v = 1$，否则$x_v = 0$。
    \item 证明这个舍入方案给出2-近似：$\text{ALG} \leq 2 \cdot \text{OPT}$。
\end{subproblems}
\end{problem}

\begin{problem}[概念题 $\star\star\star$]
简述Goemans-Williamson算法求解Max-Cut问题的思想：
\begin{subproblems}
    \item 写出Max-Cut的SDP松弛。
    \item 解释随机超平面舍入（random hyperplane rounding）的思想。
    \item 为什么近似比是$0.878 \approx \min_{0 \leq \theta \leq \pi} \frac{2}{\pi} \cdot \frac{\theta}{1 - \cos\theta}$？
\end{subproblems}
\end{problem}

\begin{problem}[思考题 $\star\star$]
为什么在GPU时代，用于求解大规模LP的一阶方法（如PDHG）正在复兴？请从以下角度分析：
\begin{subproblems}
    \item 单纯形法和内点法各自的瓶颈是什么？
    \item PDHG（Primal-Dual Hybrid Gradient）为什么适合GPU？
    \item 在什么规模/结构的LP问题上，GPU一阶方法更有优势？
\end{subproblems}
\end{problem}

%========================================
% 第4章
%========================================

% --- 按章节归并的补充题 ---
\subsection*{三、LP对偶手算（第3讲）}

\begin{problem}[$\star$]
写出以下LP的对偶问题，并画出可行域求解：
\[
\min \quad x_1 + 2x_2 \quad \text{s.t.} \quad x_1 + x_2 \geq 3, \quad 2x_1 + x_2 \geq 4, \quad x_1,x_2 \geq 0
\]
\end{problem}

\begin{problem}[$\star\star$]
\textbf{(对偶互补松弛手算)} 考虑LP及其对偶：
\[
\text{原: } \min \begin{pmatrix}2\\3\\4\end{pmatrix}^\top x \quad \text{s.t.} \quad \begin{pmatrix}1&1&0\\0&1&1\end{pmatrix}x = \begin{pmatrix}2\\3\end{pmatrix}, \quad x \geq 0
\]
\begin{subproblems}
    \item 写出对偶问题。
    \item 验证$x^* = (0, 2, 1)^\top$是原问题的可行解，计算目标值。
    \item 利用互补松弛条件$x_j^*(c_j - \sum_i a_{ij}y_i^*) = 0$确定对偶最优$y^*$。
    \item 验证强对偶成立。
\end{subproblems}
\end{problem}

\begin{problem}[$\star\star$]
\textbf{(最短路的LP对偶)} 在如下有向图上求从$s$到$t$的最短路：
\begin{center}
$s \xrightarrow{3} A \xrightarrow{2} t$, \quad $s \xrightarrow{6} t$, \quad $s \xrightarrow{1} B \xrightarrow{4} A$, \quad $B \xrightarrow{5} t$
\end{center}
\begin{subproblems}
    \item 将最短路写成LP：对每条边$e$定义流量变量$f_e$，写出目标函数和流守恒约束。
    \item 写出其LP对偶，并解释对偶变量$y_v$（每个节点一个）的含义（到$t$的最短距离）。
    \item 手算原问题和对偶问题的最优解。
\end{subproblems}
\end{problem}

\begin{problem}[$\star\star\star$]
\textbf{(博弈论中的LP对偶)} 考虑零和博弈的收益矩阵：
\[
A = \begin{pmatrix} 3 & 1 \\ 0 & 4 \end{pmatrix}
\]
行玩家选择混合策略$p = (p_1, p_2)$，列玩家选择$q = (q_1, q_2)$。
\begin{subproblems}
    \item 行玩家的maxmin问题为$\max_p \min_q p^\top A q$。将其写成LP。
    \item 写出该LP的对偶，并说明它恰好是列玩家的minmax问题$\min_q \max_p p^\top A q$。
    \item 求解纳什均衡混合策略$p^*, q^*$和博弈值$v^*$。
\end{subproblems}
\end{problem}

%--- 动态规划手算（对应第4讲）---
\subsection*{第3讲：LP对偶与近似算法}

\begin{problem}[$\star\star$]
\textbf{Vertex Cover: LP松弛舍入}

\begin{subproblems}
    \item 用 \texttt{networkx} 随机生成一个$n=20$的Erd\H{o}s-R\'{e}nyi图$G(n, p=0.3)$。
    \item 用CVXPY求解最小顶点覆盖的LP松弛。
    \item 实现舍入算法（$x_v \geq 0.5$则选入覆盖），验证结果是合法的顶点覆盖。
    \item 比较LP最优值、舍入解大小、（若$n$不太大可暴力搜索）真实最小覆盖大小。验证$\text{ALG} \leq 2\cdot\text{OPT}$。
\end{subproblems}
\end{problem}

%--- 第4讲 ---
\section{第4讲：算法设计范式与GPU化}

\begin{problem}[概念题 $\star$]
分治法的三个步骤是什么？写出分治法的时间复杂度递推关系的一般形式$T(n) = aT(n/b) + f(n)$，并说明Master定理的三种情况。
\end{problem}

\begin{problem}[计算题 $\star$]
使用Master定理分析以下递推的时间复杂度：
\begin{subproblems}
    \item $T(n) = 2T(n/2) + O(n)$（归并排序）
    \item $T(n) = 8T(n/2) + O(n^2)$（朴素矩阵乘法）
    \item $T(n) = 7T(n/2) + O(n^2)$（Strassen算法）
    \item $T(n) = 2T(n/2) + O(n \log n)$
\end{subproblems}
\end{problem}

\begin{problem}[概念题 $\star\star$]
动态规划需要满足两个条件：\textbf{最优子结构}和\textbf{重叠子问题}。
\begin{subproblems}
    \item 分别解释这两个条件的含义。
    \item 为什么贪心算法不需要"重叠子问题"？
    \item 动态规划和分治法的核心区别是什么？
\end{subproblems}
\end{problem}

\begin{problem}[计算题 $\star\star$]
考虑编辑距离问题：将字符串$s_1$变为$s_2$的最小操作数（插入、删除、替换各花费1）。
\begin{subproblems}
    \item 写出状态转移方程$dp[i][j]$。
    \item 计算"kitten"和"sitting"的编辑距离。
    \item 该DP算法的时间和空间复杂度分别是多少？
    \item 解释如何利用GPU的对角线并行（wavefront并行）来加速编辑距离的计算。为什么同一条反对角线上的元素可以并行计算？
\end{subproblems}
\end{problem}

\begin{problem}[证明题 $\star\star$]
\textbf{贪心算法的交换论证}。考虑活动选择问题：给定$n$个活动，每个活动有开始时间$s_i$和结束时间$f_i$。请用交换论证（exchange argument）证明：按结束时间最早贪心选择可以得到最优解。
\end{problem}

\begin{problem}[思考题 $\star\star\star$]
GPU暴力法的"逆袭"：
\begin{subproblems}
    \item 对于Two Sum问题（在数组中找两个数使得和为目标值），哈希表解法的复杂度是$O(n)$。GPU暴力解法检查所有$O(n^2)$对，但可以完全并行化。在什么规模的$n$下，GPU暴力可能更快？解释原因。
    \item 更一般地，讨论"顺序依赖性"如何影响算法的GPU可并行化程度。将分治、DP、贪心三种范式按GPU友好程度排序并解释原因。
\end{subproblems}
\end{problem}

%========================================
% 第5章
%========================================

% --- 按章节归并的补充题 ---
\subsection*{四、动态规划手算（第4讲）}

\begin{problem}[$\star$]
\textbf{(0-1背包)} 物品重量和价值分别为：
\begin{center}
\begin{tabular}{c|cccc}
物品 & 1 & 2 & 3 & 4 \\
\hline
重量 & 2 & 3 & 4 & 5 \\
价值 & 3 & 4 & 5 & 6
\end{tabular}
\end{center}
背包容量$W = 8$。
\begin{subproblems}
    \item 写出状态转移方程$dp[i][w]$。
    \item 手算填写完整的DP表格。
    \item 回溯得到最优选择方案。
    \item 该问题的LP松弛（将$x_i \in \{0,1\}$放宽为$0 \leq x_i \leq 1$）最优值是多少？与整数最优值的差距是多少？
\end{subproblems}
\end{problem}

\begin{problem}[$\star\star$]
\textbf{(矩阵链乘法)} 计算矩阵连乘$A_1 A_2 A_3 A_4$的最少乘法次数，其中矩阵维度分别为$A_1: 10\times 30$，$A_2: 30\times 5$，$A_3: 5\times 20$，$A_4: 20\times 15$。
\begin{subproblems}
    \item 写出状态转移方程$m[i][j] = \min_{i \leq k < j}\{m[i][k] + m[k+1][j] + p_{i-1}p_k p_j\}$。
    \item 填写完整的DP表格（$m[i][j]$和$s[i][j]$）。
    \item 最优的加括号方案是什么？
    \item 该DP表格中，哪些元素可以在GPU上并行计算？（按对角线分析）
\end{subproblems}
\end{problem}

\begin{problem}[$\star\star$]
\textbf{(最长公共子序列)} 求字符串 $X=\text{``ABCBDAB''}$ 和 $Y=\text{``BDCABA''}$ 的最长公共子序列（LCS）。
\begin{subproblems}
    \item 写出状态转移方程。
    \item 填写完整的DP表格。
    \item LCS的长度是多少？给出一个具体的LCS。
    \item 分析该DP的反对角线并行性：第$k$条反对角线上有多少个可并行计算的元素？
\end{subproblems}
\end{problem}

%--- 梯度下降手算（对应第5讲）---
\subsection*{第4讲：算法范式与GPU化}

\begin{problem}[$\star\star$]
\textbf{编辑距离的GPU对角线并行}

\begin{subproblems}
    \item 实现标准的$O(mn)$ CPU版本编辑距离。
    \item 实现GPU对角线并行版本：按反对角线$k=0,1,\ldots,m+n$扫描，每条反对角线内的元素用PyTorch向量化计算。
    \item 生成两个长度为5000的随机字符串。比较CPU版和GPU版的运行时间。
    \item 画出对角线填充的动态示意图（例如在$10\times 10$的小例子上，用颜色表示填充顺序）。
\end{subproblems}
\end{problem}

%--- 第5讲 ---
\section{第5讲：梯度法与预条件}

\begin{problem}[概念题 $\star$]
写出梯度下降法的迭代公式，并解释为什么负梯度方向是局部最速下降方向。
\end{problem}

\begin{problem}[计算题 $\star$]
考虑函数$f(x_1, x_2) = 2x_1^2 + x_2^2 - 2x_1x_2 + 4x_1 - 6x_2$。
\begin{subproblems}
    \item 计算梯度$\nabla f$和Hessian矩阵$H$。
    \item 验证$f$是强凸函数。
    \item 求$f$的最小值点（令$\nabla f = 0$）。
    \item 计算条件数$\kappa = \lambda_{\max}(H)/\lambda_{\min}(H)$，并解释它对梯度下降收敛速度的影响。
\end{subproblems}
\end{problem}

\begin{problem}[证明题 $\star\star$]
证明梯度下降法在$L$-光滑凸函数上的收敛速率。具体来说，设$f$是凸函数且$\nabla f$是$L$-Lipschitz的（即$\|\nabla f(x) - \nabla f(y)\| \leq L\|x - y\|$），步长$\eta = 1/L$。证明：
\[
f(x_T) - f(x^*) \leq \frac{L\|x_0 - x^*\|^2}{2T}
\]
\begin{hint}
先证明descent lemma: $f(y) \leq f(x) + \nabla f(x)^\top(y-x) + \frac{L}{2}\|y-x\|^2$。
\end{hint}
\end{problem}

\begin{problem}[计算题 $\star\star$]
\textbf{条件数的影响}。考虑二次函数$f(x) = \frac{1}{2}x^\top A x - b^\top x$，其中
\[
A = \begin{pmatrix} 100 & 0 \\ 0 & 1 \end{pmatrix}, \quad b = \begin{pmatrix} 100 \\ 1 \end{pmatrix}
\]
\begin{subproblems}
    \item 计算最优解$x^*$和条件数$\kappa$。
    \item 梯度下降从$x_0 = (0,0)^\top$出发，步长$\eta = 1/\lambda_{\max}(A) = 0.01$。写出前3步的迭代过程。
    \item 对于强凸函数，梯度下降的收敛率为$O\left(\left(\frac{\kappa-1}{\kappa+1}\right)^T\right)$。当$\kappa = 100$时，需要大约多少步使得误差缩小到初始的$1/e$？
\end{subproblems}
\end{problem}

\begin{problem}[概念题 $\star\star$]
比较以下优化方法，填表：
\begin{center}
\begin{tabular}{l|c|c|c}
\toprule
方法 & 每步复杂度 & 收敛速度 & GPU友好度 \\
\midrule
梯度下降 & & & \\
牛顿法 & & & \\
L-BFGS & & & \\
对角预条件(Adam) & & & \\
\bottomrule
\end{tabular}
\end{center}
并解释为什么在GPU时代，"一阶方法+简单预条件"往往优于"二阶方法"。
\end{problem}

\begin{problem}[计算题 $\star\star\star$]
\textbf{牛顿法}。考虑$f(x) = e^x + e^{-x}$。
\begin{subproblems}
    \item 计算$f'(x)$和$f''(x)$。
    \item 写出牛顿法迭代公式：$x_{k+1} = x_k - \frac{f'(x_k)}{f''(x_k)}$。
    \item 从$x_0 = 1$出发，计算前3步迭代的值。
    \item 解释为什么牛顿法具有局部二次收敛性（即$|x_{k+1} - x^*| \leq C|x_k - x^*|^2$）。
\end{subproblems}
\end{problem}

\begin{problem}[思考题 $\star\star$]
预条件的本质是"椭圆变圆"。
\begin{subproblems}
    \item 解释预条件梯度下降$x \leftarrow x - \eta P^{-1}\nabla f(x)$中$P$的作用。理想情况下$P$应该选什么？
    \item 为什么在实践中不直接用$P = H$（Hessian矩阵）？
    \item GPU友好的预条件谱系为：对角 $\to$ 块对角 $\to$ Kronecker $\to$ 低秩。解释这个谱系中"结构化 = 可并行"的含义。
\end{subproblems}
\end{problem}

%========================================
% 第6章
%========================================

% --- 按章节归并的补充题 ---
\subsection*{五、梯度下降与牛顿法手算（第5讲）}

\begin{problem}[$\star$]
\textbf{(Nocedal \& Wright风格)} 对函数$f(x,y) = (x-2)^2 + 4(y-3)^2$，从初始点$(0,0)$出发：
\begin{subproblems}
    \item 用精确线搜索（exact line search）的梯度下降，计算前3步的迭代点。
    \item 画出等高线和迭代轨迹的示意图。
    \item 这个函数的条件数是多少？迭代轨迹有什么特点？
\end{subproblems}
\begin{hint}
精确线搜索：$\eta_k = \arg\min_{\eta \geq 0} f(x_k - \eta \nabla f(x_k))$。对二次函数可以解析求解。
\end{hint}
\end{problem}

\begin{problem}[$\star\star$]
\textbf{(二次函数上的梯度下降)} 考虑$f(x) = \frac{1}{2}x^\top A x$，其中
\[
A = \begin{pmatrix} 5 & 1 \\ 1 & 2 \end{pmatrix}
\]
\begin{subproblems}
    \item 计算$A$的特征值和条件数。
    \item 从$x_0 = (1,1)^\top$出发，步长$\eta = 2/(\lambda_{\min}+\lambda_{\max})$（最优常数步长），手算前2步梯度下降。
    \item 牛顿法从同一点出发只需几步到达最优解？手算验证。
    \item 解释二者的差异：为什么牛顿法在二次函数上一步即达？
\end{subproblems}
\end{problem}

\begin{problem}[$\star\star$]
\textbf{(多维牛顿法)} 用牛顿法求解方程组（即求$\nabla f = 0$）：
\[
f(x,y) = x^2 + xy + y^2 - 6x - 9y + 20
\]
\begin{subproblems}
    \item 计算$\nabla f$和Hessian $H$。
    \item 牛顿法更新：$\mathbf{x}_{k+1} = \mathbf{x}_k - H^{-1}\nabla f(\mathbf{x}_k)$。从$(0,0)$出发，手算$\mathbf{x}_1$。
    \item 与直接求解$\nabla f = 0$对比，结果是否一致？为什么？
\end{subproblems}
\end{problem}

\begin{problem}[$\star\star\star$]
\textbf{(预条件梯度下降)} 考虑$f(x) = \frac{1}{2}x^\top A x - b^\top x$，其中
\[
A = \begin{pmatrix} 10 & 2 \\ 2 & 1 \end{pmatrix}, \quad b = \begin{pmatrix} 12 \\ 3 \end{pmatrix}
\]
\begin{subproblems}
    \item 计算最优解$x^* = A^{-1}b$。
    \item 不带预条件的梯度下降，从$x_0=(0,0)^\top$出发，$\eta=1/\lambda_{\max}(A)$，计算前2步。
    \item 取对角预条件$P = \text{diag}(10, 1)$，预条件梯度下降$x_{k+1} = x_k - \eta P^{-1}\nabla f(x_k)$。使用$\eta=1/\lambda_{\max}(P^{-1}A)$，计算前2步。
    \item 比较两种方法的收敛速度，用条件数解释差异。
\end{subproblems}
\end{problem}

\begin{problem}[$\star\star$]
\textbf{(Backtracking线搜索)} 对函数$f(x) = x^4 - 4x^3 + 4x^2$，从$x_0 = 3$出发做梯度下降。
\begin{subproblems}
    \item 计算$f'(3)$和下降方向$d_0 = -f'(3)$。
    \item 使用Armijo准则（backtracking）：$f(x_k + \eta d_k) \leq f(x_k) + c \eta f'(x_k)d_k$，取$c=0.1$，初始$\eta=1$，每次$\eta \leftarrow \eta/2$。确定第一步的步长$\eta_0$。
    \item 计算$x_1$。
\end{subproblems}
\end{problem}

\begin{problem}[$\star\star\star$]
\textbf{(共轭梯度法)} 用共轭梯度法求解$Ax=b$，其中
\[
A = \begin{pmatrix} 4 & 1 \\ 1 & 3 \end{pmatrix}, \quad b = \begin{pmatrix} 1 \\ 2 \end{pmatrix}
\]
从$x_0 = (2, 1)^\top$出发：
\begin{subproblems}
    \item 计算初始残差$r_0 = b - Ax_0$，初始搜索方向$p_0 = r_0$。
    \item 计算步长$\alpha_0 = r_0^\top r_0 / (p_0^\top A p_0)$，更新$x_1$。
    \item 计算$r_1$，验证$r_1 \perp r_0$。
    \item 计算$\beta_1$，$p_1$，$\alpha_1$，$x_2$。验证$x_2$是精确解。
\end{subproblems}
\end{problem}

%--- 流形优化手算（对应第6讲）---
\subsection*{第5讲：梯度法与预条件}

\begin{problem}[$\star$]
\textbf{梯度下降可视化}

对二维Rosenbrock函数$f(x,y) = (1-x)^2 + 100(y-x^2)^2$：
\begin{subproblems}
    \item 实现梯度下降（手动计算梯度），从$(-1, 1)$出发。
    \item 分别尝试步长$\eta = 10^{-4}, 10^{-3}, 5\times 10^{-3}$，各跑5000步。
    \item 在等高线图上画出三条轨迹，观察步长对收敛的影响。
    \item 实现带Backtracking线搜索（Armijo准则）的版本，同一初始点运行，对比效果。
\end{subproblems}
\end{problem}

\begin{problem}[$\star\star$]
\textbf{条件数实验}

构造$n=100$维的二次函数$f(x)=\frac{1}{2}x^\top Ax - b^\top x$，其中$A$是对角矩阵$A = \text{diag}(1, 2, \ldots, 100)$，$b = \mathbf{1}$。
\begin{subproblems}
    \item 用梯度下降（步长$\eta = 2/(\lambda_{\min}+\lambda_{\max})$），记录每步的$\|x_k - x^*\|$。需要多少步使误差降到$10^{-6}$？
    \item 用牛顿法，记录同样的收敛曲线。需要几步？
    \item 对角预条件$P = \text{diag}(A)$的预条件GD需要多少步？
    \item 改为$A = \text{diag}(1, 1, \ldots, 1, 1000)$（$\kappa=1000$），重复实验。画4条收敛曲线在一张图上。
\end{subproblems}
\end{problem}

%--- 第6讲 ---
\section{第6讲：流形优化}

\begin{problem}[概念题 $\star$]
什么是光滑流形（smooth manifold）？直观地解释"局部像欧氏空间"的含义。列举三个在优化中常见的流形及其约束条件。
\end{problem}

\begin{problem}[概念题 $\star\star$]
解释以下流形优化中的核心概念：
\begin{subproblems}
    \item 切空间（tangent space）$T_x\mathcal{M}$
    \item Riemannian梯度（Riemannian gradient）
    \item Retraction映射$R_x: T_x\mathcal{M} \to \mathcal{M}$
    \item 欧氏梯度与Riemannian梯度之间的关系
\end{subproblems}
\end{problem}

\begin{problem}[计算题 $\star\star$]
考虑单位球面$S^{n-1} = \{x \in \mathbb{R}^n : \|x\|_2 = 1\}$上的优化问题：
\[
\min_{x \in S^{n-1}} f(x) = x^\top A x
\]
其中$A$是对称矩阵。
\begin{subproblems}
    \item $S^{n-1}$上点$x$处的切空间$T_x S^{n-1}$是什么？
    \item 将欧氏梯度$\nabla f(x) = 2Ax$投影到切空间，得到Riemannian梯度。
    \item 写出球面上的retraction映射（归一化投影）。
    \item 这个优化问题的最优解是什么？（与$A$的特征值/特征向量有什么关系？）
\end{subproblems}
\end{problem}

\begin{problem}[计算题 $\star\star\star$]
\textbf{Stiefel流形}$\text{St}(n,p) = \{W \in \mathbb{R}^{n \times p} : W^\top W = I_p\}$。
\begin{subproblems}
    \item 证明$W$处的切空间为$T_W \text{St}(n,p) = \{\Delta : W^\top \Delta + \Delta^\top W = 0\}$。
    \item 给定欧氏梯度$G = \nabla f(W)$，Riemannian梯度为$\text{grad} f(W) = G - W(W^\top G + G^\top W)/2$。验证$\text{grad} f(W) \in T_W \text{St}(n,p)$。
    \item 简述QR-based retraction的计算过程：$R_W(\eta) = \text{qf}(W + \eta)$。
\end{subproblems}
\end{problem}

\begin{problem}[思考题 $\star\star$]
流形优化 vs 惩罚法：
\begin{subproblems}
    \item 用惩罚法处理正交约束$W^\top W = I$，写出目标函数$\tilde{f}(W) = f(W) + \frac{\mu}{2}\|W^\top W - I\|_F^2$的形式。这种方法有什么问题？
    \item 流形优化的优势是什么？为什么说它是"在约束空间上直接优化"？
    \item 举一个流形优化在深度学习中的应用（如正交RNN），并解释为什么正交约束可以缓解梯度消失问题。
\end{subproblems}
\end{problem}

%========================================
% 第7章
%========================================

% --- 按章节归并的补充题 ---
\subsection*{六、流形优化手算（第6讲）}

\begin{problem}[$\star\star$]
\textbf{(球面上的Rayleigh商)} 求$\min_{x \in S^2} x^\top A x$，其中$S^2$是$\mathbb{R}^3$中的单位球面，
\[
A = \begin{pmatrix} 3 & 1 & 0 \\ 1 & 2 & 1 \\ 0 & 1 & 1 \end{pmatrix}
\]
\begin{subproblems}
    \item 用拉格朗日乘子法：$L(x,\lambda) = x^\top A x - \lambda(x^\top x - 1)$。写出最优性条件。
    \item 说明最优解是$A$的最小特征值对应的特征向量。
    \item 计算$A$的特征值（提示：特征多项式为$\lambda^3 - 6\lambda^2 + 9\lambda - 2 = 0$，有一个根接近$\lambda \approx 0.298$）。
    \item 用Riemannian梯度下降从$x_0 = (1,0,0)^\top$出发，计算一步更新（取步长$\eta = 0.1$）。
\end{subproblems}
\end{problem}

\begin{problem}[$\star\star\star$]
\textbf{(Procrustes问题)} 给定$A, B \in \mathbb{R}^{n \times p}$，求Stiefel流形上的最优正交矩阵：
\[
\min_{W \in \text{St}(p,p)} \|A - BW\|_F^2
\]
\begin{subproblems}
    \item 将目标函数展开为$\|A\|_F^2 - 2\text{tr}(A^\top BW) + \|B\|_F^2$。
    \item 证明问题等价于$\max_{W^\top W = I} \text{tr}(A^\top BW)$。
    \item 设$A^\top B$的SVD为$U\Sigma V^\top$，证明最优解为$W^* = VU^\top$。
    \item 取$A^\top B = \begin{pmatrix} 3 & 1 \\ 1 & 2 \end{pmatrix}$，手算SVD并给出$W^*$。
\end{subproblems}
\end{problem}

%--- SGD与Adam手算（对应第8讲）---
\subsection*{第6讲：流形优化}

\begin{problem}[$\star\star$]
\textbf{球面上的Riemannian梯度下降}

求矩阵$A$的最小特征值和对应特征向量：$\min_{x \in S^{n-1}} x^\top A x$。
\begin{subproblems}
    \item 生成$n=10$的随机对称矩阵$A$。
    \item 实现Riemannian梯度下降：
    \begin{enumerate}[label=\arabic*.]
        \item 欧氏梯度$g = 2Ax$
        \item 投影到切空间$\tilde{g} = g - (x^\top g)x$
        \item 更新$\tilde{x} = x - \eta \tilde{g}$
        \item Retraction（归一化）$x = \tilde{x}/\|\tilde{x}\|$
    \end{enumerate}
    \item 画出每步$x^\top A x$的值，与\texttt{np.linalg.eigh}给出的最小特征值比较。
    \item 改变步长$\eta$，观察收敛速度和稳定性。
\end{subproblems}
\end{problem}

%--- 第7讲 ---
\section{第7讲：随机化算法}

\begin{problem}[概念题 $\star$]
叙述Johnson-Lindenstrauss引理：对于$n$个点$x_1, \ldots, x_n \in \mathbb{R}^d$和$0 < \varepsilon < 1$，存在映射$f: \mathbb{R}^d \to \mathbb{R}^k$（$k = O(\log n / \varepsilon^2)$），使得对所有$i, j$：
\[
(1 - \varepsilon)\|x_i - x_j\|^2 \leq \|f(x_i) - f(x_j)\|^2 \leq (1 + \varepsilon)\|x_i - x_j\|^2
\]
解释这个引理的意义：为什么投影维度$k$与原始维度$d$无关？
\end{problem}

\begin{problem}[计算题 $\star\star$]
\textbf{随机矩阵乘法}。给定$A \in \mathbb{R}^{m \times n}$和$B \in \mathbb{R}^{n \times p}$，随机采样$k$列的近似乘法为$AB \approx CR$，其中：
\begin{subproblems}
    \item 描述按列范数比例采样的策略：$P(j \text{ 被选中}) = \frac{\|A_{:,j}\|\|B_{j,:}\|}{\sum_l \|A_{:,l}\|\|B_{l,:}\|}$。
    \item 证明该估计是无偏的：$\mathbb{E}[CR] = AB$。
    \item 如果$A$是$10000\times 5000$，$B$是$5000\times 8000$，精确乘法需要多少次浮点运算？采样$k=500$列时需要多少次？
\end{subproblems}
\end{problem}

\begin{problem}[概念题 $\star\star$]
\textbf{Sketch-and-Solve}用于随机最小二乘。考虑$\min_x \|Ax - b\|_2$，其中$A \in \mathbb{R}^{n \times d}$，$n \gg d$。
\begin{subproblems}
    \item Sketch-and-Solve方法将问题转化为$\min_x \|SAx - Sb\|_2$，其中$S \in \mathbb{R}^{k \times n}$是sketching矩阵。$k$一般选多大？
    \item 列举两种常用的sketching矩阵（如高斯随机矩阵、CountSketch）。
    \item 解释为什么这种方法可以将时间从$O(nd^2)$降到$O(\text{nnz}(A) \cdot d + d^3)$。
\end{subproblems}
\end{problem}

\begin{problem}[计算题 $\star\star\star$]
\textbf{随机SVD}。给定矩阵$A \in \mathbb{R}^{m \times n}$，要求其top-$k$奇异值和奇异向量。
\begin{subproblems}
    \item 描述随机SVD算法的步骤：
    \begin{enumerate}[label=\arabic*.]
        \item 生成随机矩阵$\Omega \in \mathbb{R}^{n \times (k+p)}$
        \item 计算$Y = A\Omega$
        \item 对$Y$做QR分解$Y = QR$
        \item 计算$B = Q^\top A$
        \item 对$B$做SVD
    \end{enumerate}
    \item 解释为什么$Q$的列空间近似$A$的top-$k$左奇异向量张成的空间。
    \item 该算法的时间复杂度是多少？与精确SVD（$O(mn\min(m,n))$）相比有何优势？
    \item 过采样参数$p$（一般取$5 \sim 10$）的作用是什么？
\end{subproblems}
\end{problem}

\begin{problem}[思考题 $\star\star$]
解释"随机化让不可能变可能"这一观点。以$50000 \times 50000$矩阵的top-100 SVD为例，精确SVD需要的内存和时间分别是多少？随机SVD如何解决这个问题？
\end{problem}

%========================================
% 第8章
%========================================

% --- 按章节归并的补充题 ---
\subsection*{八、随机化算法手算（第7讲）}

\begin{problem}[$\star\star$]
\textbf{(JL投影手算)} 将$\mathbb{R}^4$中的三个点$x_1=(1,0,0,0)$, $x_2=(0,1,0,0)$, $x_3=(1,1,0,0)$用随机矩阵$R \in \mathbb{R}^{2 \times 4}$投影到$\mathbb{R}^2$：
\[
R = \frac{1}{\sqrt{2}}\begin{pmatrix} 1 & 1 & -1 & 0 \\ 0 & -1 & 0 & 1 \end{pmatrix}
\]
\begin{subproblems}
    \item 计算投影后的点$Rx_1, Rx_2, Rx_3$。
    \item 计算原始空间中所有点对的距离$\|x_i - x_j\|$。
    \item 计算投影空间中所有点对的距离$\|Rx_i - Rx_j\|$。
    \item 计算每个点对的距离保持比率$\|Rx_i - Rx_j\|^2/\|x_i - x_j\|^2$。是否接近1？
\end{subproblems}
\end{problem}

\begin{problem}[$\star\star$]
\textbf{(随机SVD小例子)} 考虑矩阵
\[
A = \begin{pmatrix} 4 & 0 \\ 3 & -5 \\ 0 & 4 \end{pmatrix}
\]
取随机向量$\omega = (1, 0)^\top$，按随机SVD算法的思路求$A$的rank-1近似：
\begin{subproblems}
    \item 计算$y = A\omega$。
    \item 归一化$q = y / \|y\|$。
    \item 计算$B = q^\top A$。
    \item 近似$A \approx q B$。计算近似误差$\|A - qB\|_F$。
    \item 与$A$的精确最佳rank-1近似（由SVD给出）比较。
\end{subproblems}
\end{problem}

%--- Bellman方程与值迭代手算（对应第12讲）---
\subsection*{第7讲：随机化算法}

\begin{problem}[$\star\star$]
\textbf{随机SVD实现}

\begin{subproblems}
    \item 生成一个$1000\times 1000$的低秩（rank-10）加噪声矩阵：$A = UV^\top + 0.01 \cdot E$，其中$U,V \in \mathbb{R}^{1000\times 10}$，$E$是随机噪声。
    \item 实现随机SVD算法（约15行代码）：随机投影 $\to$ QR $\to$ 投影回 $\to$ 小SVD。
    \item 与\texttt{np.linalg.svd}比较：(i) 前10个奇异值的误差；(ii) 运行时间。
    \item 增加过采样参数$p$从0到20，画出近似误差随$p$的变化曲线。
\end{subproblems}
\end{problem}

%--- 第8讲 ---
\section{第8讲：随机优化}

\begin{problem}[概念题 $\star$]
写出随机梯度下降（SGD）的更新公式，并证明SGD的梯度估计是无偏的：$\mathbb{E}[\nabla f_i(x)] = \nabla f(x)$，其中$f(x) = \frac{1}{n}\sum_{i=1}^n f_i(x)$。
\end{problem}

\begin{problem}[计算题 $\star\star$]
\textbf{Mini-batch方差分析}。设全梯度$\nabla f(x) = \frac{1}{n}\sum_{i=1}^n \nabla f_i(x)$，单个样本梯度的方差为$\text{Var}[\nabla f_i(x)] = \sigma^2$。
\begin{subproblems}
    \item 计算mini-batch（大小为$b$）梯度估计的方差$\text{Var}\left[\frac{1}{b}\sum_{j=1}^b \nabla f_{i_j}(x)\right]$。
    \item 当$b$增大$k$倍时，方差缩小多少倍？计算量增加多少倍？
    \item 为什么说存在一个"最优batch size"的trade-off？
\end{subproblems}
\end{problem}

\begin{problem}[证明题 $\star\star$]
证明SGD在凸、$L$-光滑函数上，使用步长$\eta_t = \frac{c}{\sqrt{T}}$（常数步长缩减方案）的收敛速率为$O(1/\sqrt{T})$。
\begin{hint}
使用SGD更新的递推：$\mathbb{E}\|x_{t+1} - x^*\|^2 \leq \mathbb{E}\|x_t - x^*\|^2 - 2\eta_t(f(x_t) - f(x^*)) + \eta_t^2 \sigma^2$。
\end{hint}
\end{problem}

\begin{problem}[计算题 $\star\star$]
\textbf{Adam优化器}。写出Adam的完整更新公式：
\begin{subproblems}
    \item 一阶矩估计$m_t$和二阶矩估计$v_t$的更新公式。
    \item 偏差修正$\hat{m}_t$和$\hat{v}_t$。
    \item 参数更新公式。
    \item 解释为什么Adam可以看作"对角预条件+动量"。预条件矩阵是什么？
    \item Adam的默认超参数$\beta_1, \beta_2, \epsilon$通常取什么值？
\end{subproblems}
\end{problem}

\begin{problem}[概念题 $\star\star$]
比较以下优化器的特点：
\begin{subproblems}
    \item SGD vs SGD+Momentum：动量的物理直觉是什么？
    \item Momentum vs Nesterov Momentum：Nesterov的"先预测再求梯度"有什么好处？
    \item Adam vs AdamW：AdamW中的weight decay和L2正则化有什么区别？
\end{subproblems}
\end{problem}

\begin{problem}[思考题 $\star\star\star$]
\textbf{大Batch训练}。
\begin{subproblems}
    \item 解释为什么简单地增大batch size并线性增大学习率可能导致训练不稳定。
    \item LARS（Layer-wise Adaptive Rate Scaling）的核心思想是什么？
    \item 从GPU利用率的角度，解释为什么batch=32时GPU利用率只有30\%，而batch=512时可以达到90\%。
    \item "小batch浪费GPU"这个观点与"大batch泛化性差"之间如何权衡？
\end{subproblems}
\end{problem}

%========================================
% 第9章
%========================================

% --- 按章节归并的补充题 ---
\subsection*{七、SGD与Adam手算（第8讲）}

\begin{problem}[$\star\star$]
\textbf{(SGD手算)} 考虑线性回归$f(w) = \frac{1}{3}\sum_{i=1}^{3} (w^\top x_i - y_i)^2$，数据为：
\[
(x_1, y_1) = ((1,1)^\top, 3), \quad (x_2, y_2) = ((1,-1)^\top, 1), \quad (x_3, y_3) = ((2,0)^\top, 4)
\]
\begin{subproblems}
    \item 计算全梯度$\nabla f(w)$，在$w_0 = (0,0)^\top$处求值。
    \item SGD每轮随机选一个样本。假设选取顺序为$i=2, 3, 1$，步长$\eta = 0.1$。手算SGD的前3步：$w_1, w_2, w_3$。
    \item 比较$w_3$与全梯度下降3步（$\eta = 0.1$）的结果$w_3^{GD}$。
    \item 计算单样本梯度的方差$\text{Var}[\nabla f_i(w_0)]$。
\end{subproblems}
\end{problem}

\begin{problem}[$\star\star\star$]
\textbf{(Adam优化器手算)} 对一维函数$f(w) = w^4$，从$w_0 = 2$出发，使用Adam优化器（$\beta_1 = 0.9$, $\beta_2 = 0.999$, $\epsilon = 10^{-8}$, $\eta = 0.01$）。
\begin{subproblems}
    \item 初始化$m_0 = 0$, $v_0 = 0$。
    \item 第1步：计算$g_1 = f'(w_0)$，更新$m_1, v_1, \hat{m}_1, \hat{v}_1$，最后更新$w_1$。写出每个中间值。
    \item 第2步：计算$g_2 = f'(w_1)$，重复上述步骤得到$w_2$。
    \item 与普通梯度下降（$\eta = 0.01$）的前2步比较。Adam的步长有何自适应效果？
\end{subproblems}
\end{problem}

\begin{problem}[$\star\star$]
\textbf{(Momentum手算)} 对$f(x) = 5x_1^2 + x_2^2$，从$x_0 = (5, 1)^\top$出发，$v_0 = (0,0)^\top$。使用带动量的SGD：
\[
v_{t+1} = \mu v_t - \eta \nabla f(x_t), \quad x_{t+1} = x_t + v_{t+1}
\]
取$\mu = 0.9$, $\eta = 0.05$。
\begin{subproblems}
    \item 手算前3步$(v_1, x_1), (v_2, x_2), (v_3, x_3)$。
    \item 不带动量（$\mu = 0$, $\eta = 0.05$）的前3步$x_1^{GD}, x_2^{GD}, x_3^{GD}$。
    \item 比较两种方法在$x_1$方向（病态方向）上的收敛速度。
\end{subproblems}
\end{problem}

\begin{problem}[$\star\star$]
\textbf{(Nesterov Momentum手算)} 对同样的$f(x) = 5x_1^2 + x_2^2$，从$x_0 = (5, 1)^\top$出发。Nesterov动量更新为：
\[
v_{t+1} = \mu v_t - \eta \nabla f(x_t + \mu v_t), \quad x_{t+1} = x_t + v_{t+1}
\]
取$\mu = 0.9$, $\eta = 0.05$, $v_0 = (0,0)^\top$。
\begin{subproblems}
    \item 手算前2步。注意区别：先做"lookahead" $x_t + \mu v_t$，再计算梯度。
    \item 与普通Momentum对比，前2步的$x$有什么差异？
\end{subproblems}
\end{problem}

%--- 随机化算法手算（对应第7讲）---
\subsection*{第8讲：随机优化}

\begin{problem}[$\star\star$]
\textbf{优化器赛马：SGD/Momentum/Adam}

用PyTorch在MNIST上训练一个简单的两层全连接网络（784-128-10）：
\begin{subproblems}
    \item 分别使用SGD（$\eta=0.01$）、SGD+Momentum（$\mu=0.9$）、Adam（默认参数），各训练10个epoch。
    \item 画出三条训练loss曲线和三条测试accuracy曲线。
    \item 对SGD，分别尝试batch size $\in \{16, 64, 256, 1024\}$，画出"每epoch时间 vs 收敛速度"的trade-off。
    \item 实现学习率warmup（前5\%步线性增长），观察对大batch ($\geq 512$) 训练稳定性的影响。
\end{subproblems}
\end{problem}

\begin{problem}[$\star\star\star$]
\textbf{从零实现Adam}

不使用\texttt{torch.optim}，手写Adam优化器（约20行核心代码）：
\begin{subproblems}
    \item 实现完整的Adam更新：$m$和$v$的指数移动平均、偏差修正、参数更新。
    \item 在Rosenbrock函数上与\texttt{torch.optim.Adam}的结果对比，验证实现正确。
    \item 去掉偏差修正，观察训练前期的行为差异。
    \item 将$\beta_2$改为0.9（远小于默认0.999），观察收敛行为的变化。
\end{subproblems}
\end{problem}

%--- 第9讲 ---
\section{第9讲：深度学习优化理论}

\begin{problem}[概念题 $\star$]
解释"过参数化悖论"：为什么参数量远大于数据量的深度网络不会严重过拟合？经典学习理论对此有什么预测？实际情况如何？
\end{problem}

\begin{problem}[证明题 $\star\star$]
\textbf{隐式正则化}。考虑欠定线性回归$\min_w \|Xw - y\|^2$，其中$X \in \mathbb{R}^{n \times d}$，$n < d$（有无穷多解）。证明：从$w_0 = 0$出发的梯度下降收敛到最小范数解：
\[
w^* = X^\top(XX^\top)^{-1}y = \arg\min_{Xw=y} \|w\|_2
\]
\begin{hint}
分两步：(1) 证明梯度$\nabla L = X^\top(Xw - y)$始终在$\text{row}(X)$中，因此$w_t$始终在$\text{row}(X)$中；(2) 在$\text{row}(X)$中满足$Xw = y$的解唯一，且就是最小范数解。
\end{hint}
\end{problem}

\begin{problem}[概念题 $\star\star$]
\textbf{高维鞍点}。解释为什么在高维非凸优化中，大多数临界点是鞍点而非局部极小值。
\begin{subproblems}
    \item Bray-Dean定理说了什么？
    \item 什么是Wigner半圆律？临界点处Hessian的特征值服从什么分布？
    \item 为什么临界点是局部极小的概率约为$2^{-d}$（$d$为维度）？
    \item Choromanska等人的结果如何将此与神经网络联系起来？
\end{subproblems}
\end{problem}

\begin{problem}[概念题 $\star\star$]
\textbf{逃离鞍点}。
\begin{subproblems}
    \item 定义strict saddle point。
    \item Perturbed GD如何帮助逃离鞍点？其基本思想是什么？
    \item 在负曲率方向上，扰动以什么速率增长？
    \item 为什么SGD本身的噪声可以起到类似的扰动作用？
\end{subproblems}
\end{problem}

\begin{problem}[计算题 $\star\star\star$]
\textbf{Neural Tangent Kernel (NTK)}。考虑参数为$\theta$的神经网络$f(x;\theta)$。
\begin{subproblems}
    \item 写出NTK的定义：$K(x, x') = \nabla_\theta f(x;\theta)^\top \nabla_\theta f(x';\theta)$。
    \item 在梯度流（gradient flow）下，证明网络输出的演化满足：$\frac{df}{dt} = -K(f - y)$，其中$f = (f(x_1;\theta), \ldots, f(x_n;\theta))$。
    \item 在无限宽极限下，为什么$K$趋于固定（与$\theta$无关）？这如何保证全局收敛？
    \item NTK理论的主要局限是什么？什么是"lazy training"？
\end{subproblems}
\end{problem}

\begin{problem}[思考题 $\star\star$]
训练过程中，权重矩阵的奇异值分布如何变化？
\begin{subproblems}
    \item 初始化时，权重矩阵的奇异值分布接近什么？（Marchenko-Pastur分布）
    \item 训练后观察到什么现象？（heavy tail）
    \item 为什么说"偏离MP分布越多，学得越好"？
\end{subproblems}
\end{problem}

%========================================
% 第10章
%========================================

% --- 按章节归并的补充题 ---
\subsection*{第9讲：深度学习优化理论}

\begin{problem}[$\star\star$]
\textbf{隐式正则化验证}

\begin{subproblems}
    \item 生成欠定线性回归数据：$X \in \mathbb{R}^{20 \times 100}$（随机高斯），$y = Xw_{\text{true}} + \epsilon$。
    \item 从$w_0 = 0$出发做梯度下降直到$\|Xw - y\| < 10^{-6}$。记录$\|w_T\|_2$。
    \item 计算最小范数解$w^* = X^\top(XX^\top)^{-1}y$，验证$\|w_T - w^*\|$是否很小。
    \item 改为从$w_0 = \mathbf{10}$（大初始值）出发，GD收敛到的解$\|w\|_2$是否还是最小范数？为什么？
\end{subproblems}
\end{problem}

\begin{problem}[$\star\star\star$]
\textbf{Hessian谱的随机Lanczos估计}

\begin{subproblems}
    \item 在MNIST上训练一个小的全连接网络（784-64-10）。
    \item 用PyTorch的\texttt{torch.autograd.functional.hvp}（Hessian-vector product）实现随机Lanczos算法，估计Hessian的前20个特征值。
    \item 画出Hessian特征值的直方图（随机采样100个特征值）。是否观察到"bulk + outliers"结构？
    \item 在训练过程中（每个epoch结束时）重复上述估计，画出最大特征值随训练的变化曲线。
\end{subproblems}
\begin{hint}
Lanczos算法核心：用$Hv$（而非完整$H$）做三对角化。每步只需一次Hessian-vector product，复杂度$O(d)$而非$O(d^2)$。
\end{hint}
\end{problem}

%--- 第10讲 ---
\section{第10讲：变分法与最优控制}

\begin{problem}[概念题 $\star$]
什么是泛函（functional）？给出泛函的定义，并举两个优化泛函的经典例子。
\end{problem}

\begin{problem}[计算题 $\star\star$]
\textbf{Euler-Lagrange方程}。考虑泛函$J[y] = \int_a^b F(x, y, y')\,dx$，其中$y(a)$和$y(b)$给定。
\begin{subproblems}
    \item 写出Euler-Lagrange方程：$\frac{\partial F}{\partial y} - \frac{d}{dx}\frac{\partial F}{\partial y'} = 0$。
    \item 对泛函$J[y] = \int_0^1 (y'^2 + y^2)\,dx$，$y(0) = 0$, $y(1) = 1$，写出对应的Euler-Lagrange方程。
    \item 求解上述方程，找到使$J$取极值的函数$y(x)$。
\end{subproblems}
\end{problem}

\begin{problem}[证明题 $\star\star$]
\textbf{最大熵原理}。在约束$\int p(x)\,dx = 1$和$\mathbb{E}[x] = \mu$, $\text{Var}(x) = \sigma^2$下，最大化微分熵$H[p] = -\int p(x)\log p(x)\,dx$。
\begin{subproblems}
    \item 写出带Lagrange乘子的泛函。
    \item 利用变分法导出最优分布$p^*(x)$。
    \item 验证$p^*(x)$是高斯分布$\mathcal{N}(\mu, \sigma^2)$。
\end{subproblems}
\end{problem}

\begin{problem}[计算题 $\star\star$]
\textbf{KL散度与变分推断}。
\begin{subproblems}
    \item 写出KL散度$D_{KL}(q\|p) = \int q(x)\log\frac{q(x)}{p(x)}\,dx$的定义，并证明$D_{KL}(q\|p) \geq 0$（利用Jensen不等式）。
    \item 在变分推断中，我们用$q(z)$近似后验$p(z|x)$。证明最小化$D_{KL}(q\|p(\cdot|x))$等价于最大化ELBO（Evidence Lower Bound）：
    \[
    \text{ELBO}(q) = \mathbb{E}_{q}[\log p(x,z)] - \mathbb{E}_q[\log q(z)]
    \]
\end{subproblems}
\end{problem}

\begin{problem}[概念题 $\star\star\star$]
\textbf{最优控制与HJB方程}。
\begin{subproblems}
    \item 写出连续时间最优控制问题的一般形式：$\min_{u(\cdot)} \int_0^T L(x,u)\,dt + \Phi(x(T))$，s.t. $\dot{x} = f(x,u)$。
    \item 写出Hamilton-Jacobi-Bellman（HJB）方程。
    \item 解释HJB方程与Bellman方程（RL中的离散版本）之间的联系。
    \item 伴随方法（adjoint method）如何与反向传播联系？在Neural ODE中，为什么adjoint method可以将显存从$O(L)$降到$O(1)$？
\end{subproblems}
\end{problem}

%========================================
% 第11章
%========================================

% --- 按章节归并的补充题 ---
\subsection*{十三、变分法与最优控制手算（第10讲）}

\begin{problem}[$\star\star$]
\textbf{(最速降线)} 求使物体在重力下从$(0,0)$滑到$(x_1, y_1)$所用时间最短的曲线。时间泛函为：
\[
T[y] = \int_0^{x_1} \sqrt{\frac{1 + y'^2}{2gy}} \, dx
\]
\begin{subproblems}
    \item 写出被积函数$F(y, y', x) = \sqrt{(1+y'^2)/(2gy)}$。
    \item 由于$F$不显含$x$，利用Beltrami恒等式$F - y'\frac{\partial F}{\partial y'} = C$简化Euler-Lagrange方程。
    \item 证明解为摆线的参数方程：$x = a(\theta - \sin\theta)$, $y = a(1 - \cos\theta)$。
\end{subproblems}
\end{problem}

\begin{problem}[$\star\star$]
\textbf{(简单最优控制)} 考虑一维系统$\dot{x} = u$，$x(0) = 1$，要在$T=1$时到达$x(1) = 0$，同时最小化能量$J = \int_0^1 \frac{1}{2}u^2 \, dt$。
\begin{subproblems}
    \item 写出Hamiltonian $H(x, p, u) = \frac{1}{2}u^2 + p \cdot u$，其中$p$是伴随变量。
    \item 由$\partial H / \partial u = 0$求最优控制$u^* = u^*(p)$。
    \item 写出伴随方程$\dot{p} = -\partial H / \partial x$并求解。
    \item 利用边界条件确定$u^*(t)$和最优轨迹$x^*(t)$。
    \item 计算最小能量$J^*$。
\end{subproblems}
\end{problem}

\begin{problem}[$\star\star\star$]
\textbf{(LQR手算)} 线性二次调节器：$\dot{x} = ax + bu$，目标$\min \int_0^\infty (qx^2 + ru^2)\,dt$。
取$a=1, b=1, q=1, r=1$。
\begin{subproblems}
    \item 写出代数Riccati方程$0 = 2aP - P^2 b^2/r + q$（一维形式）。
    \item 求解$P$。
    \item 最优控制为$u^* = -(b/r)Px$。写出闭环系统$\dot{x} = (a - b^2P/r)x$。
    \item 验证闭环系统是稳定的（特征值为负）。
    \item 从$x(0) = 1$出发，写出最优轨迹$x^*(t)$和$u^*(t)$。计算最优代价$J^* = Px(0)^2$。
\end{subproblems}
\end{problem}

%--- Bandit与在线学习手算（对应第11讲）---
\subsection*{第10讲：变分法与最优控制}

\begin{problem}[$\star\star$]
\textbf{Neural ODE迷你实现}

用PyTorch实现一个最简单的Neural ODE：
\begin{subproblems}
    \item 定义一个2维ODE动力学$\dot{z} = f_\theta(z)$，其中$f_\theta$是两层MLP（2-32-2）。
    \item 用\texttt{torchdiffeq.odeint}（或手写4阶Runge-Kutta）从$z(0)$积分到$z(1)$。
    \item 目标：学习将同心圆上的点分类。生成数据：内圈标签0，外圈标签1。
    \item 训练后，画出$z(t)$在$t \in [0,1]$上的连续变换过程（流场图）。
\end{subproblems}
\end{problem}

%--- 第11讲 ---
\section{第11讲：在线优化与Regret框架}

\begin{problem}[概念题 $\star$]
描述在线学习（online learning）的基本框架：每轮做什么决策？收到什么反馈？目标是什么？
\end{problem}

\begin{problem}[概念题 $\star$]
写出Regret的定义：
\[
\text{Regret}_T = \sum_{t=1}^T \ell_t(x_t) - \min_{x \in \mathcal{X}} \sum_{t=1}^T \ell_t(x)
\]
解释为什么我们追求$\text{Regret}_T = o(T)$（即$\text{Regret}_T / T \to 0$），以及这意味着什么。
\end{problem}

\begin{problem}[证明题 $\star\star$]
\textbf{Multiplicative Weights Update}。对于$N$个专家的在线决策问题，MWU算法如下：
\begin{enumerate}[label=\arabic*.]
    \item 初始化权重$w_1^i = 1$，$\forall i$
    \item 每轮$t$：选择分布$p_t^i = w_t^i / \sum_j w_t^j$
    \item 收到损失$\ell_t^i$后更新$w_{t+1}^i = w_t^i (1 - \eta \ell_t^i)$
\end{enumerate}
证明MWU的Regret界为$O(\sqrt{T \log N})$。
\begin{hint}
使用势函数$\Phi_t = \sum_i w_t^i$，分析$\Phi_t$的变化。
\end{hint}
\end{problem}

\begin{problem}[计算题 $\star\star$]
\textbf{Multi-Armed Bandit}。考虑$K=3$个臂的Bandit问题，真实期望奖励分别为$\mu_1=0.3$，$\mu_2=0.7$，$\mu_3=0.5$。
\begin{subproblems}
    \item 写出UCB1算法的选择准则：$a_t = \arg\max_i \left[\hat{\mu}_i + \sqrt{\frac{2\ln t}{n_i}}\right]$，并解释"乐观面对不确定"的含义。
    \item 假设前3轮分别拉了臂1、2、3，奖励分别为0、1、0。计算第4轮每个臂的UCB值，决定拉哪个臂。
    \item UCB1的Regret上界是$O(\sqrt{KT \log T})$。这与完全信息（在线学习）的$O(\sqrt{T \log K})$相比如何？信息缺失的代价体现在哪里？
\end{subproblems}
\end{problem}

\begin{problem}[概念题 $\star\star$]
叙述Hoeffding不等式，并解释它如何用于构造置信区间。如果对一个均值为$\mu$的$[0,1]$随机变量采样$n$次，样本均值$\hat{\mu}$满足$|\hat{\mu} - \mu| \leq \epsilon$的概率至少为多少？
\end{problem}

\begin{problem}[思考题 $\star\star\star$]
课程提出"Regret是统一语言"，从在线学习到Bandit到MDP，Regret框架如何统一它们？
\begin{subproblems}
    \item 在线学习的Regret与Bandit的Regret在定义上有什么区别？
    \item MDP的Regret如何定义？
    \item 这三者的Regret下界分别是多少？信息量从多到少如何影响Regret？
\end{subproblems}
\end{problem}

%========================================
% 第12章
%========================================

% --- 按章节归并的补充题 ---
\subsection*{十四、Bandit与在线学习手算（第11讲）}

\begin{problem}[$\star\star$]
\textbf{(Multiplicative Weights手算)} 3个专家，学习率$\eta = 0.5$。初始权重$w_1 = (1, 1, 1)$。

\medskip
\begin{tabular}{c|ccc}
轮次$t$ & 专家1损失 & 专家2损失 & 专家3损失 \\
\hline
1 & 1 & 0 & 0.5 \\
2 & 0 & 1 & 0.5 \\
3 & 0.5 & 0.5 & 1
\end{tabular}

\begin{subproblems}
    \item 手算每轮的权重$w_t$和选择分布$p_t$（归一化权重）。
    \item 计算算法的累积损失$\sum_t \langle p_t, \ell_t \rangle$。
    \item 计算最优单个专家的累积损失$\min_i \sum_t \ell_t^i$。
    \item Regret是多少？
\end{subproblems}
\end{problem}

\begin{problem}[$\star\star$]
\textbf{(Thompson Sampling手算)} 2个臂的Bernoulli Bandit。先验$\text{Beta}(1,1)$（均匀分布）。

前4轮的结果为：
\begin{center}
$t=1$: 拉臂1, 奖励=1; \quad $t=2$: 拉臂2, 奖励=0; \quad $t=3$: 拉臂1, 奖励=1; \quad $t=4$: 拉臂2, 奖励=1
\end{center}

\begin{subproblems}
    \item 每轮结束后写出每个臂的后验分布$\text{Beta}(\alpha_i, \beta_i)$。
    \item 计算每个臂在第5轮的后验均值$\alpha_i/(\alpha_i + \beta_i)$。
    \item Thompson Sampling在第5轮如何做决策？（从每个臂的后验中采样，选较大的）
    \item 与UCB的决策比较（计算第5轮的UCB值）。
\end{subproblems}
\end{problem}

\begin{problem}[$\star\star\star$]
\textbf{(Online Gradient Descent手算)} 在线凸优化：决策集$\mathcal{X} = [-1, 1]$，对手选择损失函数。
\begin{align*}
&t=1: \ell_1(x) = (x-0.5)^2, \quad t=2: \ell_2(x) = (x+0.3)^2, \\
&t=3: \ell_3(x) = (x-0.8)^2, \quad t=4: \ell_4(x) = (x+0.1)^2
\end{align*}
OGD更新：$\tilde{x}_{t+1} = x_t - \eta \nabla \ell_t(x_t)$，$x_{t+1} = \Pi_\mathcal{X}(\tilde{x}_{t+1})$（投影回$[-1,1]$）。

取$x_1 = 0$, $\eta = 0.5$。
\begin{subproblems}
    \item 手算4轮的$x_t$和$\ell_t(x_t)$。
    \item 计算算法的总损失$\sum_t \ell_t(x_t)$。
    \item 找到事后最优$x^* = \arg\min_{x \in [-1,1]} \sum_t \ell_t(x)$（可对总损失求导）。
    \item 计算Regret。
\end{subproblems}
\end{problem}

%--- 综合计算题 ---
\subsection*{十五、综合计算题}

\begin{problem}[$\star\star\star$]
\textbf{(从对偶到算法)} 考虑岭回归问题：
\[
\min_w \frac{1}{2}\|Xw - y\|^2 + \frac{\lambda}{2}\|w\|^2
\]
其中$X \in \mathbb{R}^{n \times d}$，$n < d$。
\begin{subproblems}
    \item 直接求解：$w^* = (X^\top X + \lambda I)^{-1}X^\top y$。需要求逆一个$d \times d$矩阵。
    \item 利用矩阵恒等式证明：$w^* = X^\top(XX^\top + \lambda I)^{-1}y$（对偶形式），只需求逆一个$n \times n$矩阵。
    \item 取$X = \begin{pmatrix} 1 & 2 & 0 \\ 0 & 1 & 1 \end{pmatrix}$, $y = (3, 2)^\top$, $\lambda = 1$，分别用两种形式手算$w^*$并验证结果一致。
    \item 当$d \gg n$时，哪种形式更高效？这体现了什么"对偶"思想？
\end{subproblems}
\end{problem}

\begin{problem}[$\star\star\star$]
\textbf{(完整优化流水线)} 考虑二维Rosenbrock函数$f(x,y) = (1-x)^2 + 100(y-x^2)^2$，最优解为$(1,1)$。
\begin{subproblems}
    \item 计算$(0,0)$处的梯度和Hessian。
    \item 从$(0,0)$出发，计算一步梯度下降（$\eta = 0.001$）。
    \item 从$(0,0)$出发，计算一步牛顿法。Hessian是否正定？如果不是，如何修正？
    \item 解释Rosenbrock函数为什么是经典的"病态"测试函数（从条件数角度分析）。
\end{subproblems}
\end{problem}

\begin{problem}[$\star\star\star$]
\textbf{(Fisher信息矩阵手算)} 考虑二元分类的logistic回归。模型$p(y=1|x;\theta) = \sigma(\theta^\top x)$，其中$\sigma(z) = 1/(1+e^{-z})$。
\begin{subproblems}
    \item 证明$\nabla_\theta \log p(y|x;\theta) = (y - \sigma(\theta^\top x))x$。
    \item 写出Fisher信息矩阵$F(\theta) = \mathbb{E}[\nabla \log p \cdot (\nabla \log p)^\top]$。
    \item 给定单个数据点$x = (1, 1)^\top$，$\theta = (0, 0)^\top$，计算$F(\theta)$。
    \item 比较普通梯度$\nabla_\theta J$和自然梯度$F^{-1}\nabla_\theta J$在该点的区别。
\end{subproblems}
\end{problem}

\begin{problem}[$\star\star$]
\textbf{(NTK手算)} 考虑最简单的两层线性网络$f(x; w_1, w_2) = w_2 \cdot w_1 \cdot x$，其中$x, w_1, w_2 \in \mathbb{R}$。
\begin{subproblems}
    \item 计算$\nabla_{(w_1,w_2)} f = (w_2 x, \; w_1 x)$。
    \item 给定两个输入$x_1 = 1, x_2 = -1$，写出NTK矩阵$K_{ij} = \nabla f(x_i)^\top \nabla f(x_j)$。
    \item 取$w_1 = w_2 = 1$，计算具体的$K$矩阵。
    \item 若目标$y = (1, -1)$，写出梯度流方程$\frac{d}{dt}\begin{pmatrix}f(x_1)\\f(x_2)\end{pmatrix} = -K \begin{pmatrix}f(x_1)-y_1\\f(x_2)-y_2\end{pmatrix}$，并分析$K$的特征值对学习速度的影响。
\end{subproblems}
\end{problem}

%========================================
%========================================
% 补充：小规模编程题
%========================================
%========================================
\newpage
\subsection*{第11讲：在线优化与Bandit}

\begin{problem}[$\star\star$]
\textbf{Bandit算法对比赛}

实现$\varepsilon$-greedy、UCB1和Thompson Sampling三种算法：
\begin{subproblems}
    \item 设置10臂Bernoulli Bandit，真实概率$\mu = (0.1, 0.2, \ldots, 1.0)$。
    \item 每种算法独立运行$T=10000$步，重复50次取平均。
    \item 画出三种算法的"累积Regret vs 时间步"曲线。
    \item 画出每种算法在$t=100, 1000, 10000$时各臂的拉取次数柱状图。
    \item 对$\varepsilon$-greedy，调节$\varepsilon \in \{0.01, 0.05, 0.1, 0.3\}$，观察对Regret的影响。
\end{subproblems}
\end{problem}

%--- 第12讲 ---
\section{第12讲：MDP与Tabular RL理论}

\begin{problem}[概念题 $\star$]
定义MDP的五元组$(S, A, P, R, \gamma)$，并分别解释每个元素的含义。写出有限时间（episodic）和无限时间（discounted）两种设定下的目标函数。
\end{problem}

\begin{problem}[计算题 $\star\star$]
\textbf{Bellman方程}。考虑一个简单的GridWorld：$2\times 2$的网格，状态$\{s_1, s_2, s_3, s_4\}$（左上、右上、左下、右下），动作$\{$上、下、左、右$\}$。$s_4$是终点（奖励+1），每步奖励为$-0.04$，折扣因子$\gamma = 0.9$。
\begin{subproblems}
    \item 写出最优Bellman方程$V^*(s) = \max_a \left[R(s,a) + \gamma \sum_{s'} P(s'|s,a)V^*(s')\right]$。
    \item 使用值迭代（Value Iteration），从$V_0(s) = 0$出发，计算$V_1$和$V_2$。
    \item 值迭代为什么一定收敛？与Bellman算子的收缩映射性质有什么关系？
\end{subproblems}
\end{problem}

\begin{problem}[证明题 $\star\star$]
证明Bellman最优算子$\mathcal{T}$是$\gamma$-收缩映射，即对任意值函数$V_1, V_2$：
\[
\|\mathcal{T}V_1 - \mathcal{T}V_2\|_\infty \leq \gamma \|V_1 - V_2\|_\infty
\]
由此说明值迭代的收敛性（Banach不动点定理）。
\end{problem}

\begin{problem}[概念题 $\star\star$]
比较值迭代（Value Iteration）和策略迭代（Policy Iteration）：
\begin{subproblems}
    \item 分别描述两种算法的步骤。
    \item 哪种算法的每步计算量更大？哪种收敛更快（迭代次数）？
    \item 策略迭代为什么能保证在有限步内收敛？
\end{subproblems}
\end{problem}

\begin{problem}[概念题 $\star\star\star$]
\textbf{UCB-VI算法}。
\begin{subproblems}
    \item UCB-VI是如何将Bandit中的"乐观原则"推广到MDP的？
    \item 其Regret上界为$\tilde{O}(\sqrt{H^3 S A T})$，其中$H$是episode长度，$S$是状态数，$A$是动作数，$T$是总步数。解释每个因子的含义。
    \item Regret下界为$\Omega(\sqrt{HSAT})$。与上界相比，差距在哪里？
    \item 什么是$(\varepsilon, \delta)$-PAC？它与Regret有什么关系？
\end{subproblems}
\end{problem}

\begin{problem}[思考题 $\star\star$]
GPU上的值迭代：
\begin{subproblems}
    \item 为什么值迭代适合GPU并行？哪些计算可以并行？
    \item 当状态空间$|S| = 100000$时，为什么GPU版本比CPU版本快50倍？
    \item 表格方法的根本限制是什么？为什么需要函数逼近？
\end{subproblems}
\end{problem}

%========================================
% 第13章
%========================================

% --- 按章节归并的补充题 ---
\subsection*{九、Bellman方程与值迭代手算（第12讲）}

\begin{problem}[$\star\star$]
\textbf{(Sutton \& Barto风格)} 考虑$3 \times 3$ GridWorld：
\begin{center}
\begin{tabular}{|c|c|c|}
\hline
$s_1$ & $s_2$ & $s_3$ \\
\hline
$s_4$ & $s_5$ & $s_6$ \\
\hline
$s_7$ & $s_8$ & $s_9$ \\
\hline
\end{tabular}
\end{center}
$s_9$是目标状态（到达得$+10$奖励），其余每步$-1$奖励。动作：上、下、左、右（确定性）。碰墙不动。$\gamma = 0.9$。
\begin{subproblems}
    \item 从$V_0 = \mathbf{0}$出发做值迭代。手算$V_1(s)$（所有9个状态）。
    \item 手算$V_2(s)$（所有9个状态）。
    \item 从$V_2$中提取贪心策略$\pi_2(s) = \arg\max_a [R + \gamma V_2(s')]$。这个策略合理吗？
    \item 如果所有状态同时更新（同步），与按某个顺序逐个更新（异步/Gauss-Seidel）相比，结果有什么不同？
\end{subproblems}
\end{problem}

\begin{problem}[$\star\star$]
\textbf{(策略评估手算)} 给定策略$\pi$：在$s_1$选右，在$s_2$选右，在$s_3$选下。使用线性方程组法（而非迭代法）求$V^\pi$。
\begin{center}
$s_1 \xrightarrow[\text{r=1}]{\text{右}} s_2 \xrightarrow[\text{r=2}]{\text{右}} s_3 \xrightarrow[\text{r=0}]{\text{下}} \text{终止}$
\end{center}
$\gamma = 0.9$。
\begin{subproblems}
    \item 写出Bellman期望方程组：$V^\pi(s) = R(s,\pi(s)) + \gamma V^\pi(s')$。
    \item 求解线性方程组得到$V^\pi(s_1), V^\pi(s_2), V^\pi(s_3)$。
\end{subproblems}
\end{problem}

\begin{problem}[$\star\star\star$]
\textbf{(策略迭代完整手算)} 考虑MDP：$S = \{s_1, s_2\}$，$A = \{a, b\}$，$\gamma = 0.5$。
\[
P(s_1|s_1,a) = 0.7, \quad P(s_2|s_1,a) = 0.3, \quad R(s_1,a) = 5
\]
\[
P(s_1|s_1,b) = 0.3, \quad P(s_2|s_1,b) = 0.7, \quad R(s_1,b) = 10
\]
\[
P(s_1|s_2,a) = 0.5, \quad P(s_2|s_2,a) = 0.5, \quad R(s_2,a) = -1
\]
\[
P(s_1|s_2,b) = 0.9, \quad P(s_2|s_2,b) = 0.1, \quad R(s_2,b) = 2
\]
\begin{subproblems}
    \item \textbf{策略评估}：从初始策略$\pi_0(s_1) = a, \pi_0(s_2) = a$出发，求解$V^{\pi_0}(s_1), V^{\pi_0}(s_2)$（线性方程组）。
    \item \textbf{策略改进}：对每个状态，检查每个动作的$Q$值：$Q^{\pi_0}(s,a') = R(s,a') + \gamma \sum_{s'} P(s'|s,a')V^{\pi_0}(s')$。确定改进后的策略$\pi_1$。
    \item 对$\pi_1$重复策略评估和策略改进，直到策略不再变化。
    \item 与值迭代（从$V_0 = 0$出发迭代5轮）的结果比较。
\end{subproblems}
\end{problem}

%--- TD学习与策略梯度手算（对应第13讲）---
\subsection*{第12讲：MDP与Tabular RL}

\begin{problem}[$\star\star$]
\textbf{值迭代与策略迭代：GridWorld}

实现一个$10\times 10$的GridWorld（确定性转移，四个角落有奖励，中间有障碍物），$\gamma = 0.95$。
\begin{subproblems}
    \item 实现值迭代，记录每轮的$\|V_{k+1} - V_k\|_\infty$，画收敛曲线。
    \item 实现策略迭代（策略评估用线性方程组直接求解），记录策略改变次数。
    \item 比较两种方法到达$\|V - V^*\|_\infty < 10^{-4}$所需的迭代轮数和总时间。
    \item 用\texttt{matplotlib}可视化最终的$V^*$热力图和$\pi^*$箭头图。
\end{subproblems}
\end{problem}

\begin{problem}[$\star\star\star$]
\textbf{GPU并行值迭代}

\begin{subproblems}
    \item 随机生成一个$|S|=50000$, $|A|=4$的MDP（稀疏转移矩阵，每个$(s,a)$只转移到5个后继状态）。
    \item 用NumPy实现CPU版值迭代。
    \item 用PyTorch（\texttt{cuda} tensor）实现GPU版值迭代（核心：$V(s) = \max_a [R(s,a) + \gamma \sum_{s'} P(s'|s,a)V(s')]$向量化）。
    \item 对比CPU和GPU版本运行100轮迭代的时间。画出加速比随$|S|$变化的曲线（$|S| \in \{1000, 5000, 10000, 50000\}$）。
\end{subproblems}
\end{problem}

%--- 第13讲 ---
\section{第13讲：深度RL与策略优化}

\begin{problem}[概念题 $\star$]
解释DQN（Deep Q-Network）的核心思想：
\begin{subproblems}
    \item 为什么需要经验回放（Experience Replay）？
    \item 为什么需要目标网络（Target Network）？
    \item DQN使用的损失函数是什么？
\end{subproblems}
\end{problem}

\begin{problem}[证明题 $\star\star$]
\textbf{策略梯度定理}。证明：
\[
\nabla_\theta J(\theta) = \mathbb{E}_{\tau \sim \pi_\theta}\left[\sum_{t=0}^{T} \nabla_\theta \log \pi_\theta(a_t|s_t) \cdot Q^{\pi_\theta}(s_t, a_t)\right]
\]
\begin{hint}
从$J(\theta) = \mathbb{E}_{\tau \sim \pi_\theta}[R(\tau)]$出发，利用$\nabla_\theta \pi_\theta = \pi_\theta \nabla_\theta \log \pi_\theta$。
\end{hint}
\end{problem}

\begin{problem}[计算题 $\star\star$]
\textbf{REINFORCE与方差缩减}。
\begin{subproblems}
    \item 写出REINFORCE算法的梯度估计：$\hat{g} = \sum_t \nabla_\theta \log \pi_\theta(a_t|s_t) \cdot G_t$，其中$G_t$是什么？
    \item 引入baseline $b(s_t)$后，梯度估计变为$\hat{g} = \sum_t \nabla_\theta \log \pi_\theta(a_t|s_t) \cdot (G_t - b(s_t))$。证明任意baseline $b$不改变梯度的期望（即仍然无偏）。
    \item 最优baseline $b^*(s)$是什么？为什么通常取$b(s) \approx V^{\pi}(s)$？
\end{subproblems}
\end{problem}

\begin{problem}[概念题 $\star\star$]
\textbf{Actor-Critic方法}。
\begin{subproblems}
    \item Actor和Critic分别学什么？
    \item 什么是TD误差（Temporal Difference error）$\delta_t = r_t + \gamma V(s_{t+1}) - V(s_t)$？它在Actor-Critic中起什么作用？
    \item 与REINFORCE相比，Actor-Critic的优势和劣势分别是什么？
\end{subproblems}
\end{problem}

\begin{problem}[计算题 $\star\star\star$]
\textbf{自然梯度与Fisher信息矩阵}。
\begin{subproblems}
    \item 写出Fisher信息矩阵的定义：$F(\theta) = \mathbb{E}_{\pi_\theta}\left[\nabla_\theta \log \pi_\theta(a|s) \nabla_\theta \log \pi_\theta(a|s)^\top\right]$。
    \item 自然梯度更新为$\theta \leftarrow \theta - \eta F^{-1} \nabla_\theta J$。解释为什么这比普通梯度更好（参数化不变性）。
    \item Fisher信息矩阵与KL散度有什么关系？（提示：$F$是KL散度的Hessian）
    \item 这与课程第5讲的"预条件"暗线有什么联系？
\end{subproblems}
\end{problem}

\begin{problem}[概念题 $\star\star$]
\textbf{TRPO}。
\begin{subproblems}
    \item 写出TRPO的优化问题：$\max_\theta L(\theta)$ s.t. $D_{KL}(\pi_{\theta_\text{old}} \| \pi_\theta) \leq \delta$。
    \item 为什么需要信任域约束？如果不加约束会怎样？
    \item TRPO使用什么方法近似求解这个约束优化问题？
\end{subproblems}
\end{problem}

%========================================
% 第14章
%========================================

% --- 按章节归并的补充题 ---
\subsection*{十、TD学习与策略梯度手算（第13讲）}

\begin{problem}[$\star\star$]
\textbf{(TD(0)手算)} 给定以下经验轨迹（在同一策略$\pi$下）和$\gamma = 1$：
\begin{align*}
&\text{轨迹1: } s_1 \xrightarrow{r=0} s_2 \xrightarrow{r=1} \text{终止} \\
&\text{轨迹2: } s_1 \xrightarrow{r=0} s_2 \xrightarrow{r=1} \text{终止} \\
&\text{轨迹3: } s_2 \xrightarrow{r=0} \text{终止}
\end{align*}
初始$V(s_1) = V(s_2) = 0$，学习率$\alpha = 0.1$。
\begin{subproblems}
    \item 依次用每条轨迹的每一步做TD(0)更新：$V(s) \leftarrow V(s) + \alpha[r + \gamma V(s') - V(s)]$。写出每次更新后$V(s_1)$和$V(s_2)$的值。
    \item 蒙特卡洛方法（first-visit）得到的$V(s_1)$和$V(s_2)$估计是多少？
    \item 两种方法的结果有什么区别？解释原因。
\end{subproblems}
\end{problem}

\begin{problem}[$\star\star\star$]
\textbf{(REINFORCE手算)} 考虑一个1状态的Bandit问题：$|S|=1$，$A=\{0,1\}$。策略参数化为softmax：
\[
\pi_\theta(a=1) = \frac{e^\theta}{1 + e^\theta} = \sigma(\theta), \quad \pi_\theta(a=0) = 1 - \sigma(\theta)
\]
奖励$R(a=1) = 1$，$R(a=0) = 0$。取$\theta_0 = 0$（初始时两个动作等概率），学习率$\eta = 0.5$。
\begin{subproblems}
    \item 计算$\nabla_\theta \log \pi_\theta(a=1)$和$\nabla_\theta \log \pi_\theta(a=0)$。
    \item 计算策略梯度的期望$\mathbb{E}_{a \sim \pi_\theta}[\nabla_\theta \log \pi_\theta(a) \cdot R(a)]$在$\theta_0 = 0$处的值。
    \item 假设实际采样到$a=1$（奖励1），REINFORCE更新后$\theta_1 = ?$
    \item 假设接下来采样到$a=0$（奖励0），更新$\theta_2 = ?$
    \item 引入baseline $b = 0.5$（期望奖励），重新计算(c)(d)，比较方差是否减小。
\end{subproblems}
\end{problem}

%--- PPO手算（对应第14讲）---
\subsection*{第13讲：深度RL与策略优化}

\begin{problem}[$\star\star$]
\textbf{REINFORCE解CartPole}

用PyTorch从零实现REINFORCE算法（不使用RL库）：
\begin{subproblems}
    \item 策略网络：两层MLP（4-64-2），输出softmax概率。
    \item 实现采样轨迹、计算折扣回报$G_t$、策略梯度更新的完整循环。
    \item 在\texttt{gymnasium} CartPole-v1上训练，画出每episode的总奖励曲线（取50个episode的滑动平均）。
    \item 加入baseline $b = \bar{G}$（该batch的平均回报），比较有无baseline的收敛速度和方差。
\end{subproblems}
\begin{hint}
核心梯度：\texttt{loss = -(log\_prob * (G - baseline)).mean()}，然后\texttt{loss.backward()}。
\end{hint}
\end{problem}

%--- 第14讲 ---
\section{第14讲：PPO——约束优化的巅峰}

\begin{problem}[概念题 $\star$]
PPO有两个变体：PPO-Penalty和PPO-Clip。分别写出它们的目标函数。
\end{problem}

\begin{problem}[计算题 $\star\star$]
\textbf{PPO-Clip}。定义概率比$r_t(\theta) = \frac{\pi_\theta(a_t|s_t)}{\pi_{\theta_\text{old}}(a_t|s_t)}$。
\begin{subproblems}
    \item PPO-Clip的目标函数为
    \[
    L^{\text{CLIP}}(\theta) = \mathbb{E}_t\left[\min\left(r_t(\theta)\hat{A}_t, \;\text{clip}(r_t(\theta), 1-\epsilon, 1+\epsilon)\hat{A}_t\right)\right]
    \]
    其中$\hat{A}_t$是优势函数估计。解释clipping操作的直觉：当$\hat{A}_t > 0$和$\hat{A}_t < 0$时，clipping分别限制了什么？
    \item 取$\epsilon = 0.2$。如果$r_t = 1.5$，$\hat{A}_t = 2$，计算clipped目标值。
    \item 如果$r_t = 0.6$，$\hat{A}_t = -1$，计算clipped目标值。
    \item 为什么说PPO-Clip近似实现了信任域约束？
\end{subproblems}
\end{problem}

\begin{problem}[概念题 $\star\star$]
\textbf{GAE（Generalized Advantage Estimation）}。
\begin{subproblems}
    \item 写出$\lambda$-return的定义和GAE的公式：$\hat{A}_t^{\text{GAE}(\gamma,\lambda)} = \sum_{l=0}^{\infty}(\gamma\lambda)^l \delta_{t+l}$。
    \item $\lambda = 0$和$\lambda = 1$分别对应什么？各自的偏差和方差如何？
    \item 为什么GAE是偏差-方差trade-off的有效工具？
\end{subproblems}
\end{problem}

\begin{problem}[思考题 $\star\star$]
从课程的"对偶"暗线角度理解PPO的演进：
\begin{subproblems}
    \item TRPO：$\max L(\theta)$ s.t. $\text{KL} \leq \delta$（硬约束）。
    \item PPO-Penalty：$\max L(\theta) - \beta \cdot \text{KL}$（软惩罚）。这与Lagrangian对偶的关系是什么？
    \item PPO-Clip：用clipping近似信任域。这从第2讲的KKT到第14讲的PPO，"约束$\to$对偶$\to$一阶"的演进说明了什么？
\end{subproblems}
\end{problem}

\begin{problem}[概念题 $\star\star\star$]
\textbf{RLHF中的PPO}。
\begin{subproblems}
    \item 描述RLHF的三个阶段：SFT $\to$ 奖励模型训练 $\to$ PPO微调。
    \item 在PPO微调阶段，目标函数是什么？为什么需要加入KL惩罚项$\beta \cdot D_{KL}(\pi_\theta \| \pi_{\text{ref}})$？
    \item PPO在RLHF中面临哪些工程挑战（如reward hacking、训练不稳定）？
\end{subproblems}
\end{problem}

%========================================
% 第15章
%========================================

% --- 按章节归并的补充题 ---
\subsection*{十一、PPO相关手算（第14讲）}

\begin{problem}[$\star\star$]
\textbf{(重要性采样手算)} 旧策略$\pi_{\text{old}}$和新策略$\pi_\theta$在某状态$s$下的动作概率如下：
\begin{center}
\begin{tabular}{c|ccc}
动作 & $a_1$ & $a_2$ & $a_3$ \\
\hline
$\pi_{\text{old}}(a|s)$ & 0.5 & 0.3 & 0.2 \\
$\pi_\theta(a|s)$ & 0.6 & 0.3 & 0.1
\end{tabular}
\end{center}
优势估计$\hat{A}(s,a_1) = 2$, $\hat{A}(s,a_2) = -1$, $\hat{A}(s,a_3) = 0$。
\begin{subproblems}
    \item 计算每个动作的概率比$r(a) = \pi_\theta(a|s)/\pi_{\text{old}}(a|s)$。
    \item 计算PPO-Clip目标（$\epsilon = 0.2$）在每个动作上的值。
    \item 计算PPO-Clip目标的期望$L^{\text{CLIP}}$（在$\pi_{\text{old}}$下采样）。
    \item 计算$D_{KL}(\pi_{\text{old}} \| \pi_\theta)$。
\end{subproblems}
\end{problem}

\begin{problem}[$\star\star$]
\textbf{(GAE手算)} 给定一段轨迹的TD误差序列：$\delta_0 = 3$, $\delta_1 = -1$, $\delta_2 = 2$, $\delta_3 = 0$。$\gamma = 0.99$。
\begin{subproblems}
    \item 计算$\lambda = 0$时的GAE：$\hat{A}_t^{(0)} = \delta_t$。
    \item 计算$\lambda = 1$时的GAE：$\hat{A}_t^{(1)} = \sum_{l=0}^{\infty} \gamma^l \delta_{t+l}$（取$t=0$）。
    \item 计算$\lambda = 0.95$时的GAE：$\hat{A}_0^{(\lambda)} = \sum_{l=0}^{3} (\gamma\lambda)^l \delta_l$（截断到已知项）。
    \item 对比三个$\hat{A}_0$的值，解释$\lambda$如何控制偏差-方差权衡。
\end{subproblems}
\end{problem}

%--- ADMM手算（对应第15讲）---
\subsection*{第14讲：PPO}

\begin{problem}[$\star\star\star$]
\textbf{PPO-Clip最小实现}

在CartPole上实现一个最小版PPO（约100行核心代码）：
\begin{subproblems}
    \item 实现Actor（策略网络）和Critic（价值网络），共享底层或分开均可。
    \item 实现数据采集：运行$N$步收集$(s, a, r, s', \log\pi_{\text{old}})$。
    \item 实现GAE计算（$\lambda=0.95$, $\gamma=0.99$）。
    \item 实现PPO-Clip更新：对采集的数据做$K=4$个epoch的minibatch更新。
    \item 画出训练曲线，与纯REINFORCE对比收敛速度和稳定性。
    \item （附加）输出训练过程中$r_t(\theta)$的分布直方图，验证clipping在起作用。
\end{subproblems}
\end{problem}

%--- 第15讲 ---
\section{第15讲：分布式优化}

\begin{problem}[概念题 $\star$]
解释分布式优化中数据并行（Data Parallelism）和模型并行（Model Parallelism）的区别。各自适用于什么场景？
\end{problem}

\begin{problem}[计算题 $\star\star$]
\textbf{ADMM}。考虑问题$\min f(x) + g(z)$ s.t. $Ax + Bz = c$。ADMM的迭代步骤为：
\begin{align*}
x^{k+1} &= \arg\min_x \left(f(x) + \frac{\rho}{2}\|Ax + Bz^k - c + u^k\|_2^2\right) \\
z^{k+1} &= \arg\min_z \left(g(z) + \frac{\rho}{2}\|Ax^{k+1} + Bz - c + u^k\|_2^2\right) \\
u^{k+1} &= u^k + Ax^{k+1} + Bz^{k+1} - c
\end{align*}
\begin{subproblems}
    \item 用ADMM求解LASSO问题：$\min \frac{1}{2}\|Ax - b\|_2^2 + \lambda\|z\|_1$ s.t. $x = z$。写出$x$-更新和$z$-更新的显式公式。
    \item $z$-更新涉及的近端算子（proximal operator）$\text{prox}_{\lambda\|\cdot\|_1}(v) = ?$
    \item 解释ADMM中对偶上升（dual ascent）的思想。$u$的更新有什么含义？
\end{subproblems}
\end{problem}

\begin{problem}[概念题 $\star\star$]
\textbf{数据并行SGD}。
\begin{subproblems}
    \item 描述同步SGD（Synchronous SGD）的过程：每个worker计算本地梯度，然后如何聚合？
    \item 异步SGD与同步SGD相比有什么优缺点？
    \item 什么是"掉队者问题"（straggler problem）？如何缓解？
\end{subproblems}
\end{problem}

\begin{problem}[计算题 $\star\star$]
\textbf{All-Reduce}。$K$个GPU需要汇总各自的梯度向量（每个长度为$n$）。
\begin{subproblems}
    \item 朴素方法（所有GPU发送到一个master，再广播）的通信量是多少？
    \item Ring All-Reduce算法的通信量是多少？解释其工作原理。
    \item 为什么Ring All-Reduce对GPU集群更友好？
\end{subproblems}
\end{problem}

\begin{problem}[概念题 $\star\star$]
\textbf{梯度压缩}。
\begin{subproblems}
    \item 解释Top-K梯度稀疏化的思想：只传输最大的$K$个梯度分量。
    \item 量化压缩（如1-bit SGD）如何工作？其核心假设是什么？
    \item 梯度压缩引入的误差如何影响收敛？如何用误差反馈（error feedback）缓解？
\end{subproblems}
\end{problem}

\begin{problem}[思考题 $\star\star\star$]
\textbf{联邦学习}。
\begin{subproblems}
    \item FedAvg算法的步骤是什么？与数据并行SGD有什么区别？
    \item 联邦学习面临的主要挑战有哪些？（如非i.i.d.数据、通信效率、隐私）
    \item 从优化的角度，为什么非i.i.d.数据分布会影响FedAvg的收敛性？
\end{subproblems}
\end{problem}

%========================================
% 综合题
%========================================
\newpage

% --- 按章节归并的补充题 ---
\subsection*{十二、ADMM与分布式优化手算（第15讲）}

\begin{problem}[$\star\star$]
\textbf{(ADMM求解LASSO)} 用ADMM求解：$\min_x \frac{1}{2}\|Ax - b\|^2 + \lambda \|x\|_1$。取
\[
A = \begin{pmatrix} 1 & 2 \\ 3 & 4 \end{pmatrix}, \quad b = \begin{pmatrix} 1 \\ 2 \end{pmatrix}, \quad \lambda = 0.5, \quad \rho = 1
\]
从$z^0 = (0,0)^\top$, $u^0 = (0,0)^\top$出发。
\begin{subproblems}
    \item $x$-更新：$x^{k+1} = (A^\top A + \rho I)^{-1}(A^\top b + \rho(z^k - u^k))$。计算$(A^\top A + \rho I)^{-1}$和$x^1$。
    \item $z$-更新：$z^{k+1} = S_{\lambda/\rho}(x^{k+1} + u^k)$，其中$S_\kappa(v)_i = \text{sign}(v_i)\max(|v_i|-\kappa, 0)$是软阈值。计算$z^1$。
    \item $u$-更新：$u^{k+1} = u^k + x^{k+1} - z^{k+1}$。计算$u^1$。
    \item 再做一轮迭代，计算$(x^2, z^2, u^2)$。
\end{subproblems}
\end{problem}

\begin{problem}[$\star\star$]
\textbf{(近端算子手算)} 计算以下函数的近端算子$\text{prox}_{\lambda f}(v) = \arg\min_x \left(f(x) + \frac{1}{2\lambda}\|x - v\|^2\right)$：
\begin{subproblems}
    \item $f(x) = \|x\|_1$（$\ell_1$范数），$v = (3, -0.3, 0.8)^\top$，$\lambda = 0.5$。
    \item $f(x) = I_C(x)$（凸集$C = \{x : \|x\| \leq 1\}$的指示函数），$v = (3, 4)^\top$。
    \item $f(x) = \frac{1}{2}x^\top Q x$（$Q$正定），一般$v$。
    \item $f(x) = \|x\|_2$（$\ell_2$范数，非平方），$v = (3, 4)^\top$，$\lambda = 2$。
\end{subproblems}
\end{problem}

\begin{problem}[$\star\star$]
\textbf{(Ring All-Reduce手算)} 3个GPU，每个GPU有梯度向量$g^{(k)} \in \mathbb{R}^3$：
\[
g^{(1)} = (1, 2, 3), \quad g^{(2)} = (4, 5, 6), \quad g^{(3)} = (7, 8, 9)
\]
按Ring All-Reduce算法：
\begin{subproblems}
    \item 将每个向量分为3个chunk（每个1个元素）。
    \item 执行Reduce-Scatter阶段（2轮），写出每轮每个GPU上每个chunk的值。
    \item 执行All-Gather阶段（2轮），写出最终每个GPU上的完整聚合结果。
    \item 验证每个GPU都得到了$g^{(1)} + g^{(2)} + g^{(3)} = (12, 15, 18)$。
    \item 总通信量是多少（发送数据量）？与朴素的"发到中心再广播"比较。
\end{subproblems}
\end{problem}

%--- 变分法手算（对应第10讲）---
\subsection*{第15讲：分布式优化}

\begin{problem}[$\star\star$]
\textbf{ADMM求解LASSO}

\begin{subproblems}
    \item 生成稀疏信号恢复问题：真实$w \in \mathbb{R}^{200}$只有10个非零分量，$X \in \mathbb{R}^{50 \times 200}$，$y = Xw + \epsilon$。
    \item 实现ADMM求解$\min \frac{1}{2}\|Xw-y\|^2 + \lambda\|w\|_1$。
    \item 画出ADMM的原始残差$\|x^k - z^k\|$和对偶残差$\|\rho(z^k - z^{k-1})\|$随迭代的变化。
    \item 改变$\rho \in \{0.1, 1, 10, 100\}$，比较收敛速度。
    \item 与\texttt{sklearn.linear\_model.Lasso}的结果对比（恢复的稀疏模式是否一致）。
\end{subproblems}
\end{problem}

\begin{problem}[$\star\star$]
\textbf{模拟分布式SGD}

在单机上模拟$K$个worker的数据并行SGD（无需多GPU）：
\begin{subproblems}
    \item 将MNIST训练集均匀分为$K=4$份，分配给4个"worker"。
    \item 实现同步SGD：每个worker在本地数据上计算梯度，取平均后更新。
    \item 实现"local SGD"（FedAvg风格）：每个worker本地跑$\tau$步，然后同步模型参数。
    \item 比较$\tau = 1$（等价于同步SGD）、$\tau = 5$、$\tau = 20$的训练曲线。
    \item 模拟non-i.i.d.数据划分（每个worker只有2-3个类别的数据），观察FedAvg的收敛变化。
\end{subproblems}
\end{problem}

%========================================
%========================================
% 补充：经典深度手算题（高阶）
%========================================
%========================================
\newpage
\section*{附：综合题（跨章节）}
\addcontentsline{toc}{section}{附：综合题（跨章节）}
\setcounter{problem}{0}
\renewcommand{\theproblem}{\arabic{problem}}

\begin{problem}[综合 $\star\star\star$]
\textbf{五条暗线追踪}。课程贯穿五条暗线：对偶、Regret、预条件、随机化、流形。请各用一段话（100字以内）总结每条暗线从第一次出现到最后一次出现的演进脉络。
\end{problem}

\begin{problem}[综合 $\star\star\star$]
\textbf{GPU范式转变}。
\begin{subproblems}
    \item 列出至少4个"经典思维 vs GPU思维"的对比，并各举课程中的一个具体例子。
    \item 在你理解中，"IO感知 $>$ FLOPS"意味着什么？用FlashAttention为例说明。
    \item 课程认为"并行度 $>$ 复杂度"。这个观点有什么局限性？在什么情况下仍然是复杂度更重要？
\end{subproblems}
\end{problem}

\begin{problem}[综合 $\star\star\star$]
\textbf{从凸到非凸到在线}。回顾课程对不同优化问题的处理方式：
\begin{subproblems}
    \item 凸优化（第2-3讲）：为什么凸性是"最重要的结构"？凸优化有哪些"礼物"（好的性质）？
    \item 非凸优化（第9讲）：在深度学习中，为什么非凸问题也能被有效求解？列出三个理论解释。
    \item 在线优化（第11-14讲）：Regret框架如何将在线学习、Bandit和RL统一起来？
    \item 从"发现结构$\to$发展方法"的角度，总结这三类问题各自利用了什么结构。
\end{subproblems}
\end{problem}

\begin{problem}[综合 $\star\star$]
\textbf{从理论到实践的差距}。课程中多次强调理论最优不等于实践最优。请各举一个例子说明：
\begin{subproblems}
    \item 理论复杂度低但实践中慢的算法。
    \item 理论复杂度高但在GPU上很快的算法。
    \item 理论上有保证但实践中不使用的方法。
    \item 理论上没有保证但实践中效果很好的方法。
\end{subproblems}
\end{problem}

%========================================
%========================================
% 补充：经典手算计算题
%========================================
%========================================
\section*{补充：经典深度手算题（高阶）}
\addcontentsline{toc}{section}{补充：经典深度手算题（高阶）}
\setcounter{problem}{0}
\renewcommand{\theproblem}{\arabic{problem}}

\bigskip
\begin{center}
\textit{以下题目计算量较大、推导链较长，建议留2--3小时专注完成。\\
每题均有明确的数值答案，可自行验证。}
\end{center}

%--------------------------------------------------
\subsection*{一、凸优化深度手算}
%--------------------------------------------------

\begin{problem}[$\star\star\star$ \textbf{Markowitz投资组合完整推导}]
三种资产，期望收益$\mu = (0.10, 0.15, 0.20)^\top$，协方差矩阵
\[
\Sigma = \begin{pmatrix} 0.04 & 0.006 & 0.002 \\ 0.006 & 0.09 & 0.018 \\ 0.002 & 0.018 & 0.16 \end{pmatrix}
\]
考虑$\min_w \frac{1}{2}w^\top \Sigma w$ \; s.t. $\mu^\top w \geq 0.15$, $\mathbf{1}^\top w = 1$, $w \geq 0$。
\begin{subproblems}
    \item 先忽略$w \geq 0$约束。写出仅含等式/不等式约束$\mu^\top w \geq 0.15$和$\mathbf{1}^\top w = 1$的Lagrangian $L(w,\lambda,\nu)$。
    \item 对$w$求导令为零：$\Sigma w = \lambda \mu + \nu \mathbf{1}$，解得$w^* = \Sigma^{-1}(\lambda \mu + \nu \mathbf{1})$。先计算$\Sigma^{-1}$（手算$3\times 3$逆矩阵）。
    \item 代入两个约束$\mu^\top w = 0.15$和$\mathbf{1}^\top w = 1$，得到关于$(\lambda, \nu)$的$2\times 2$线性方程组。定义
    \[
    a = \mu^\top \Sigma^{-1}\mu, \quad b = \mu^\top \Sigma^{-1}\mathbf{1}, \quad c = \mathbf{1}^\top\Sigma^{-1}\mathbf{1}
    \]
    计算$a,b,c$的数值，解出$\lambda^*, \nu^*$。
    \item 代回求$w^*$。检查$w^*$的各分量是否$\geq 0$。如果某个$w_i^* < 0$，则$w \geq 0$约束被激活——需要将该资产剔除（$w_i = 0$），在剩余资产上重新求解。
    \item 计算最优组合的期望收益和标准差$(\mu^\top w^*, \sqrt{w^{*\top}\Sigma w^*})$。
    \item 改变目标收益$r_{\min}$在$[0.10, 0.20]$之间取5个值，重复上述过程，在$(\sigma, \mu)$平面上画出有效前沿的草图。
\end{subproblems}
\end{problem}

\begin{problem}[$\star\star\star$ \textbf{Lagrange对偶的完整推导：支持向量机}]
数据点：
\begin{center}
\begin{tabular}{c|cccccc}
$i$ & 1 & 2 & 3 & 4 & 5 & 6 \\
\hline
$x_i$ & $(3,1)$ & $(3,-1)$ & $(6,1)$ & $(1,0)$ & $(0,1)$ & $(0,-1)$ \\
$y_i$ & $+1$ & $+1$ & $+1$ & $-1$ & $-1$ & $-1$
\end{tabular}
\end{center}
\begin{subproblems}
    \item 写出软间隔SVM（$C=10$）的原问题：
    \[
    \min_{w,b,\xi} \frac{1}{2}\|w\|^2 + C\sum_i \xi_i \quad \text{s.t.} \quad y_i(w^\top x_i + b) \geq 1-\xi_i, \quad \xi_i \geq 0
    \]
    \item 写出完整的Lagrangian和KKT条件（含$\alpha_i$和$\beta_i$两组乘子）。
    \item 消去$w, b, \xi$，推导出对偶问题：
    \[
    \max_\alpha \sum_i \alpha_i - \frac{1}{2}\sum_{i,j}\alpha_i\alpha_j y_i y_j x_i^\top x_j \quad \text{s.t.} \quad 0 \leq \alpha_i \leq C, \quad \sum_i \alpha_i y_i = 0
    \]
    \item 计算$6\times 6$的矩阵$Q_{ij} = y_i y_j x_i^\top x_j$（手算每个元素）。
    \item 利用互补松弛和KKT条件分析：猜测支持向量集合（$\alpha_i > 0$的点），列出方程组求解$\alpha^*$。
    \item 由$w^* = \sum_i \alpha_i y_i x_i$和$b^*$的计算公式，得到最终分类超平面。画出数据点、超平面和间隔带的示意图。
\end{subproblems}
\end{problem}

\begin{problem}[$\star\star\star$ \textbf{内点法一步迭代}]
考虑LP：$\min c^\top x$ s.t. $Ax=b, x \geq 0$，其中
\[
c = \begin{pmatrix}1\\2\end{pmatrix}, \quad A = \begin{pmatrix}1&1\end{pmatrix}, \quad b = (1)
\]
（即$\min x_1 + 2x_2$ s.t. $x_1+x_2=1, x \geq 0$）。

使用对数障碍法，障碍参数$t > 0$：
\[
\min_x \; t \cdot c^\top x - \sum_i \ln x_i \quad \text{s.t.} \quad Ax = b
\]
\begin{subproblems}
    \item 写出该问题的KKT系统（对$x$求导并引入等式约束乘子$\nu$）。
    \item 取$t=1$，手解KKT系统，得到分析中心$x^*(t=1)$。
    \item 取$t=10$，手解KKT系统，得到$x^*(t=10)$。
    \item 取$t=100$，手解$x^*(t=100)$。
    \item 随$t \to \infty$，$x^*(t)$趋向哪里？验证它趋向LP的最优解。
    \item 画出中心路径$\{x^*(t) : t > 0\}$在$(x_1, x_2)$平面上的轨迹。
\end{subproblems}
\end{problem}

%--------------------------------------------------
\subsection*{二、梯度法与预条件深度手算}
%--------------------------------------------------

\begin{problem}[$\star\star\star$ \textbf{L-BFGS两步迭代完整推导}]
考虑$f(x,y) = 5x^2 + y^2 - xy - 11x + 5$。

L-BFGS使用最近$m$步的$\{s_k, y_k\}$对近似$H^{-1}$，其中$s_k = x_{k+1} - x_k$，$y_k = \nabla f_{k+1} - \nabla f_k$。

\begin{subproblems}
    \item 计算$\nabla f(x,y) = (10x - y - 11, \; 2y - x)$。求最优解$x^*$（令$\nabla f = 0$）。
    \item 从$x_0 = (0,0)^\top$出发，第一步无历史信息，使用$H_0^{-1} = I$，做一步梯度下降：$x_1 = x_0 - \eta H_0^{-1}\nabla f(x_0)$。取$\eta = 0.1$，计算$x_1$。
    \item 计算$s_0 = x_1 - x_0$，$y_0 = \nabla f(x_1) - \nabla f(x_0)$，$\rho_0 = 1/(y_0^\top s_0)$。
    \item 使用L-BFGS的two-loop recursion（$m=1$）计算第二步的搜索方向$p_1$：
    \begin{enumerate}[label=\arabic*.]
        \item $q = \nabla f(x_1)$
        \item $\alpha_0 = \rho_0 s_0^\top q$, \; $q \leftarrow q - \alpha_0 y_0$
        \item $r = H_0^{(1)} q$，其中$H_0^{(1)} = \frac{s_0^\top y_0}{y_0^\top y_0}I$
        \item $\beta_0 = \rho_0 y_0^\top r$, \; $r \leftarrow r + s_0(\alpha_0 - \beta_0)$
        \item $p_1 = -r$
    \end{enumerate}
    写出每个中间变量的数值。
    \item 更新$x_2 = x_1 + p_1$（步长取1），计算$f(x_2)$。
    \item 与纯梯度下降$x_2^{GD}$对比。L-BFGS是否更快地接近了最优解？
\end{subproblems}
\end{problem}

\begin{problem}[$\star\star\star$ \textbf{最速下降法在椭圆函数上的锯齿轨迹}]
考虑$f(x) = \frac{1}{2}x^\top Ax$，其中$A = \begin{pmatrix} \kappa & 0 \\ 0 & 1 \end{pmatrix}$，$\kappa = 10$。
从$x_0 = (1,1)^\top$出发，使用精确线搜索的梯度下降。
\begin{subproblems}
    \item 对二次函数，精确线搜索步长为$\eta_k = \frac{g_k^\top g_k}{g_k^\top A g_k}$，其中$g_k = Ax_k$。验证这个公式。
    \item 手算$x_1, x_2, x_3, x_4$（4步完整迭代）。
    \item 验证相邻两步梯度正交：$g_{k+1}^\top g_k = 0$。
    \item 证明对二次函数，精确线搜索GD的误差满足
    \[
    \frac{f(x_{k+1})}{f(x_k)} \leq \left(\frac{\kappa - 1}{\kappa + 1}\right)^2
    \]
    当$\kappa = 10$时，每步最多缩小多少比例？需要多少步使误差降到$10^{-6}$？
    \item 在$(x_1, x_2)$平面上画出迭代轨迹和$f$的等高线。这就是经典的"锯齿现象"。
\end{subproblems}
\end{problem}

%--------------------------------------------------
\subsection*{三、随机化与矩阵分析深度手算}
%--------------------------------------------------

\begin{problem}[$\star\star\star$ \textbf{幂迭代与Krylov子空间}]
考虑对称矩阵
\[
A = \begin{pmatrix} 5 & 4 & 1 \\ 4 & 5 & 1 \\ 1 & 1 & 2 \end{pmatrix}
\]
\begin{subproblems}
    \item \textbf{幂迭代}：从$v_0 = (1, 0, 0)^\top$出发，做幂迭代$v_{k+1} = Av_k / \|Av_k\|$。手算$v_1, v_2, v_3$，以及对应的Rayleigh商$\lambda_k = v_k^\top A v_k$。
    \item \textbf{Lanczos}：从$q_1 = v_0 / \|v_0\| = (1,0,0)^\top$出发，做Lanczos三对角化。手算3步：
    \begin{enumerate}[label=\arabic*.]
        \item $w_1 = Aq_1$, $\alpha_1 = q_1^\top w_1$, $w_1 \leftarrow w_1 - \alpha_1 q_1$, $\beta_1 = \|w_1\|$, $q_2 = w_1/\beta_1$
        \item $w_2 = Aq_2 - \beta_1 q_1$, $\alpha_2 = q_2^\top w_2$, $w_2 \leftarrow w_2 - \alpha_2 q_2$, $\beta_2 = \|w_2\|$, $q_3 = w_2/\beta_2$
        \item $w_3 = Aq_3 - \beta_2 q_2$, $\alpha_3 = q_3^\top w_3$
    \end{enumerate}
    写出三对角矩阵$T_3 = \begin{pmatrix}\alpha_1 & \beta_1 & 0 \\ \beta_1 & \alpha_2 & \beta_2 \\ 0 & \beta_2 & \alpha_3\end{pmatrix}$。
    \item 求$T_3$的特征值（它们近似$A$的特征值）。与$A$的真实特征值（$\lambda \approx 0.443, 2.000, 9.557$）比较精度。
    \item 为什么Lanczos只需$O(n)$的Hessian-vector product就能估计极端特征值？这与深度学习中用Lanczos估计Hessian谱有什么联系？
\end{subproblems}
\end{problem}

\begin{problem}[$\star\star\star$ \textbf{Marchenko-Pastur定律验证}]
考虑$n \times p$的矩阵$X$，其元素独立同分布于$\mathcal{N}(0, 1/n)$。样本协方差矩阵为$\hat{\Sigma} = X^\top X$。当$n, p \to \infty$，$p/n \to \gamma$时，$\hat{\Sigma}$的特征值的经验分布收敛到Marchenko-Pastur分布。

取$\gamma = 0.5$（即$p = n/2$）。MP分布支撑为$[\lambda_-, \lambda_+]$，其中$\lambda_\pm = (1 \pm \sqrt{\gamma})^2$。
\begin{subproblems}
    \item 计算$\lambda_-$和$\lambda_+$的数值。
    \item MP密度函数为$\rho(\lambda) = \frac{1}{2\pi \gamma \lambda}\sqrt{(\lambda_+ - \lambda)(\lambda - \lambda_-)}$。验证$\int_{\lambda_-}^{\lambda_+}\rho(\lambda)d\lambda = 1$（利用半圆的面积公式）。
    \item 计算MP分布的均值$\mathbb{E}[\lambda] = 1$和方差$\text{Var}[\lambda] = \gamma$。
    \item 现在加入信号：$X = U_{\text{signal}} + Z$，其中$U_{\text{signal}}$是rank-1矩阵（一个强信号方向），$Z$是噪声。$\hat{\Sigma}$的特征值分布变为"bulk + 一个outlier"。BBP相变说：只有当信号强度$\sigma^2 > \sqrt{\gamma}$时，outlier才从bulk中跳出。当$\gamma = 0.5$时，临界信号强度是多少？
    \item 联系课程第9讲：为什么训练好的神经网络权重矩阵的奇异值分布偏离MP分布是"好事"？
\end{subproblems}
\end{problem}

%--------------------------------------------------
\subsection*{四、动态规划与对偶深度手算}
%--------------------------------------------------

\begin{problem}[$\star\star\star$ \textbf{最大流最小割：对偶完整手算}]
考虑如下网络（$s$为源，$t$为汇）：
\begin{center}
\begin{tabular}{lc}
$s \to A$: 容量 10 & $A \to B$: 容量 4 \\
$s \to B$: 容量 8  & $A \to t$: 容量 7 \\
$B \to A$: 容量 2  & $B \to t$: 容量 6
\end{tabular}
\end{center}
\begin{subproblems}
    \item 将最大流问题写成LP（流守恒约束+容量约束）。
    \item 手动用增广路径法（Ford-Fulkerson）求最大流。写出每条增广路径和对应的残量图。
    \item 写出LP的对偶问题，证明它等价于最小割问题。
    \item 找到最小割（即将节点分为$S$和$T = V\setminus S$两组，使得$s \in S, t \in T$，横跨割的容量最小）。
    \item 验证最大流值$=$最小割容量（强对偶/最大流最小割定理）。
\end{subproblems}
\end{problem}

\begin{problem}[$\star\star\star$ \textbf{Viterbi算法手算}]
隐马尔可夫模型（HMM）有3个隐状态$\{H_1, H_2, H_3\}$，2个观测值$\{O_1, O_2\}$。
\[
\text{转移矩阵 } A = \begin{pmatrix} 0.6 & 0.3 & 0.1 \\ 0.1 & 0.7 & 0.2 \\ 0.2 & 0.2 & 0.6 \end{pmatrix}, \quad
\text{发射矩阵 } B = \begin{pmatrix} 0.8 & 0.2 \\ 0.3 & 0.7 \\ 0.5 & 0.5 \end{pmatrix}
\]
初始分布$\pi = (0.5, 0.3, 0.2)$。观测序列为$O_1, O_2, O_1, O_1$。
\begin{subproblems}
    \item 用Viterbi算法（动态规划）求最可能的隐状态序列。写出$\delta_t(j)$（到时刻$t$隐状态为$j$的最大概率）和$\psi_t(j)$（回溯指针）的完整表格。
    \item 回溯得到最优路径。
    \item 该观测序列的似然$P(O|\lambda)$是多少？（用前向算法的$\alpha_t(j)$表格计算。）
    \item Viterbi和前向算法本质上都是DP。说明两者的状态转移方程的区别（$\max$ vs $\sum$）。
    \item 分析该DP的反对角线并行性（能否用GPU加速？瓶颈在哪？）
\end{subproblems}
\end{problem}

%--------------------------------------------------
\subsection*{五、强化学习深度手算}
%--------------------------------------------------

\begin{problem}[$\star\star\star$ \textbf{Q-learning完整手算}]
$3 \times 3$ GridWorld，起点$s_1$（左上），终点$s_9$（右下，奖励$+10$），$s_5$是"陷阱"（奖励$-5$）。其余步奖励$-1$。$\gamma = 0.9$，$\varepsilon = 0.1$。

动作空间$\{$上、下、左、右$\}$，确定性转移（碰墙不动）。

初始$Q(s,a) = 0$，学习率$\alpha = 0.5$。
给定以下经验轨迹：
\begin{align*}
s_1 \xrightarrow{\text{右},r=-1} s_2 \xrightarrow{\text{右},r=-1} s_3 \xrightarrow{\text{下},r=-1} s_6 \xrightarrow{\text{下},r=-1} s_9(+10)
\end{align*}
\begin{subproblems}
    \item 对轨迹中每一步做Q-learning更新：
    \[
    Q(s,a) \leftarrow Q(s,a) + \alpha[r + \gamma \max_{a'} Q(s',a') - Q(s,a)]
    \]
    从最后一步$s_6 \to s_9$开始（为什么从后往前更新等价于多次前向？），写出每次更新后变化的$Q$值。
    \item 再给一条轨迹：$s_1 \to s_4 \to s_5(-5) \to s_6 \to s_9(+10)$。更新所有$Q$值。
    \item 经过两条轨迹后，从$s_1$出发的$\varepsilon$-greedy策略会选什么动作？
    \item 与SARSA的更新$Q(s,a) \leftarrow Q(s,a) + \alpha[r + \gamma Q(s',a') - Q(s,a)]$对比。如果第二条轨迹中的动作也由$\varepsilon$-greedy产生（假设$s_4$处探索选了"下"落入陷阱），SARSA的更新与Q-learning有何不同？
\end{subproblems}
\end{problem}

\begin{problem}[$\star\star\star$ \textbf{策略梯度的方差分析}]
考虑一个简单的MDP：$|S|=1$（只有一个状态），$|A|=3$（三个动作），奖励为$R(a_1)=1$, $R(a_2)=2$, $R(a_3)=0$（确定性）。策略用softmax参数化：
\[
\pi_\theta(a_k) = \frac{e^{\theta_k}}{\sum_j e^{\theta_j}}, \quad \theta = (\theta_1, \theta_2, \theta_3)
\]
取$\theta = (0, 0, 0)$（均匀策略）。
\begin{subproblems}
    \item 计算$\pi_\theta(a_k) = 1/3$，$\forall k$。
    \item 对每个动作$a_k$，计算$\nabla_\theta \log \pi_\theta(a_k)$。

    提示：$\frac{\partial \log \pi_\theta(a_k)}{\partial \theta_j} = \mathbf{1}_{j=k} - \pi_\theta(a_j)$。
    \item 计算策略梯度的精确值：
    \[
    g = \mathbb{E}_{a \sim \pi}[R(a) \nabla_\theta \log \pi_\theta(a)] = \sum_k \pi_\theta(a_k) R(a_k) \nabla_\theta \log \pi_\theta(a_k)
    \]
    \item 计算梯度估计的方差$\text{Var}_{a \sim \pi}[R(a) \nabla_\theta \log \pi_\theta(a)]$（每个分量的方差）。
    \item 引入baseline $b = \bar{R} = \mathbb{E}[R(a)] = 1$，重新计算策略梯度（验证期望不变）和方差（验证方差减小）。
    \item 最优baseline $b^* = \frac{\mathbb{E}[R^2 \|\nabla \log \pi\|^2]}{\mathbb{E}[\|\nabla \log \pi\|^2]}$，计算$b^*$的值。对比$b=0$、$b=\bar{R}$和$b=b^*$三种情况的方差。
\end{subproblems}
\end{problem}

%--------------------------------------------------
\subsection*{六、变分法与最优控制深度手算}
%--------------------------------------------------

\begin{problem}[$\star\star\star$ \textbf{变分推断完整推导：高斯混合模型}]
观测$x \in \mathbb{R}$，隐变量$z \in \{1, 2\}$表示簇归属。模型：
\[
p(z=k) = \pi_k, \quad p(x|z=k) = \mathcal{N}(x; \mu_k, 1)
\]
取$\pi_1 = \pi_2 = 0.5$, $\mu_1 = -1$, $\mu_2 = 2$。观测到$x_0 = 0.5$。
\begin{subproblems}
    \item 精确后验$p(z|x_0)$：用贝叶斯公式计算$p(z=1|x_0)$和$p(z=2|x_0)$。

    （需要计算$\mathcal{N}(0.5; -1, 1)$和$\mathcal{N}(0.5; 2, 1)$的密度值。）
    \item 现在考虑完整的变分推断。假设有$n=3$个观测$\mathbf{x} = (x_1, x_2, x_3) = (-0.5, 0.3, 1.8)$，隐变量$\mathbf{z} = (z_1, z_2, z_3)$。

    Mean-field近似：$q(\mathbf{z}) = \prod_i q_i(z_i)$，其中$q_i(z_i = k) = r_{ik}$。
    \item 写出ELBO：$\text{ELBO} = \sum_i \sum_k r_{ik}[\ln \pi_k + \ln \mathcal{N}(x_i;\mu_k,1) - \ln r_{ik}]$。
    \item CAVI（坐标上升变分推断）更新：
    \[
    \ln r_{ik}^{\text{new}} \propto \ln \pi_k + \ln \mathcal{N}(x_i; \mu_k, 1) = \ln \pi_k - \frac{(x_i-\mu_k)^2}{2} + \text{const}
    \]
    然后归一化$r_{ik}$。从$r_{ik} = 0.5$出发，手算一轮CAVI更新。
    \item 计算更新后的ELBO值，验证它增大了。
\end{subproblems}
\end{problem}

\begin{problem}[$\star\star\star$ \textbf{Pontryagin极大原理：双积分器}]
系统$\ddot{q} = u$（即$\dot{x}_1 = x_2, \dot{x}_2 = u$），$x_1(0) = 1, x_2(0) = 0$。目标：在$T=2$时到达$x_1(2)=0, x_2(2)=0$，最小化$J = \int_0^2 \frac{1}{2}u^2 dt$。
\begin{subproblems}
    \item 写出Hamiltonian $H = \frac{1}{2}u^2 + p_1 x_2 + p_2 u$。
    \item 由$\partial H/\partial u = 0$得最优控制$u^* = -p_2$。
    \item 写出伴随方程：$\dot{p}_1 = -\partial H/\partial x_1 = 0$, $\dot{p}_2 = -\partial H/\partial x_2 = -p_1$。求解$p_1(t), p_2(t)$。
    \item 代入$u^* = -p_2$和状态方程，求解$x_1(t), x_2(t)$（含待定常数$p_1(0)$和$p_2(0)$）。
    \item 利用终端条件$x_1(2) = 0, x_2(2) = 0$确定$p_1(0), p_2(0)$。
    \item 写出最优控制$u^*(t)$和最优轨迹$x^*(t)$的显式表达式。计算最优代价$J^*$。
    \item 该问题的HJB方程是什么形式？验证value function $V(x,t)$满足HJB方程。
\end{subproblems}
\end{problem}

%--------------------------------------------------
\subsection*{七、综合深度手算}
%--------------------------------------------------

\begin{problem}[$\star\star\star$ \textbf{从KKT到ADMM到分布式：完整链路}]
考虑分布式LASSO：$K=2$个节点各持有部分数据$(A_1, b_1)$和$(A_2, b_2)$，共同求解：
\[
\min_w \frac{1}{2}\|A_1 w - b_1\|^2 + \frac{1}{2}\|A_2 w - b_2\|^2 + \lambda\|w\|_1
\]
取$A_1 = \begin{pmatrix}1 & 0\end{pmatrix}$, $b_1 = (2)$, $A_2 = \begin{pmatrix}0 & 1\end{pmatrix}$, $b_2 = (3)$, $\lambda = 0.5$。
\begin{subproblems}
    \item \textbf{集中式求解}：直接写KKT条件（次梯度条件），手解最优$w^*$。
    \item \textbf{共识ADMM分解}：引入局部变量$w_1, w_2$和全局变量$z$：
    \[
    \min_{w_1,w_2,z} \frac{1}{2}\|A_1 w_1-b_1\|^2 + \frac{1}{2}\|A_2 w_2-b_2\|^2 + \lambda\|z\|_1 \quad \text{s.t.} \quad w_1 = z, \; w_2 = z
    \]
    写出ADMM的$w_1$-更新、$w_2$-更新、$z$-更新（全局软阈值）。
    \item 取$\rho = 1$，从$z^0 = (0,0)^\top$, $u_1^0 = u_2^0 = (0,0)^\top$出发，手算3轮ADMM迭代。
    \item 验证ADMM的解趋近于集中式的最优解$w^*$。
    \item 在这个分解中，每个节点只需要知道什么信息？什么信息需要通信？这体现了分布式优化的什么优势？
\end{subproblems}
\end{problem}

\begin{problem}[$\star\star\star$ \textbf{Mirror Descent完整推导}]
在单纯形$\Delta_n = \{p \in \mathbb{R}^n : p_i \geq 0, \sum p_i = 1\}$上做在线优化。Mirror Descent使用负熵作为Bregman散度：$\Phi(p) = \sum_i p_i \ln p_i$。

取$n=3$（3个专家），$\eta = 0.5$。
\begin{subproblems}
    \item Bregman散度$D_\Phi(p\|q) = \sum_i p_i \ln(p_i/q_i)$（即KL散度）。验证$D_\Phi$非负且当$p=q$时为零。
    \item Mirror Descent更新：$p_{t+1} = \arg\min_{p \in \Delta_n} \left[\eta \langle \ell_t, p \rangle + D_\Phi(p \| p_t)\right]$。证明闭式解为：
    \[
    p_{t+1}^i = \frac{p_t^i \cdot e^{-\eta \ell_t^i}}{\sum_j p_t^j \cdot e^{-\eta \ell_t^j}}
    \]
    （这就是Multiplicative Weights Update！）
    \item 从$p_1 = (1/3, 1/3, 1/3)$出发，损失向量$\ell_1 = (1, 0, 0.5)$, $\ell_2 = (0, 1, 0.5)$, $\ell_3 = (0.5, 0.5, 0)$。手算$p_2, p_3, p_4$。
    \item 计算Regret $= \sum_t \langle \ell_t, p_t \rangle - \min_i \sum_t \ell_t^i$。
    \item 一般Regret界为$\frac{\ln n}{\eta} + \frac{\eta T}{2}$。取$\eta = \sqrt{2\ln n / T}$可得$O(\sqrt{T \ln n})$。当$n=3, T=3$时，理论界给出多少？与你算的实际Regret比较。
\end{subproblems}
\end{problem}

%========================================
%========================================
% 补充：开放研究思考题（从课内出发，通向前沿）
%========================================
%========================================
\newpage
\section*{补充：开放研究思考题}
\addcontentsline{toc}{section}{补充：开放研究思考题}
\setcounter{problem}{0}
\renewcommand{\theproblem}{\arabic{problem}}

\bigskip
\begin{center}
\textit{以下题目从课程知识出发，通过逐步诱导引向研究前沿。\\
没有唯一标准答案，重在推理过程。鼓励查阅文献后再回头审视自己的推导。}
\end{center}

%--------------------------------------------------
\subsection*{Scaling Law与计算最优分配}
%--------------------------------------------------

\begin{problem}[$\star\star\star$ \textbf{算力该给谁：大模型还是大Batch？}]
你有固定的GPU预算$C$（以FLOPs计）。你可以选择模型大小$N$（参数量）和训练数据量$D$（token数）。训练的总计算量近似为$C \approx 6ND$。

经验上，在充分训练后，模型的测试loss满足如下幂律关系（Scaling Law）：
\[
L(N, D) \approx \left(\frac{N_0}{N}\right)^{\alpha} + \left(\frac{D_0}{D}\right)^{\beta} + L_\infty
\]
其中$\alpha \approx 0.076$，$\beta \approx 0.095$，$L_\infty$是不可约误差，$N_0, D_0$是常数。
\begin{subproblems}
    \item \textbf{(从第8讲出发)} 回忆SGD在强凸函数上的收敛率$O(1/T)$和凸函数上的$O(1/\sqrt{T})$。对于固定模型，增加训练步数$T$（等价于增加$D$）的边际收益是递减的。用$L$关于$D$的幂律项$D^{-\beta}$解释这一点。

    \item 现在固定总预算$C = 6ND$。将$D = C/(6N)$代入，得到$L$仅关于$N$的函数$L(N) = (N_0/N)^\alpha + (6N \cdot D_0/C)^{-\beta} + L_\infty$。对$N$求导，找到最优模型大小$N^*(C)$。证明$N^* \propto C^{\beta/(\alpha+\beta)}$。

    \item 代入$\alpha=0.076, \beta=0.095$，计算$N^*$和$D^*$各随$C$的增长指数。当预算翻倍时，最优策略是模型变大多少？数据变多多少？

    \item \textbf{(关键洞察)} 一个同学说："我有很多GPU，所以我把batch size开到最大，把小模型训到极致。"请用上述分析解释为什么这在经济上是次优的——同样的FLOPs花在更大的模型上，loss会更低。给出具体的数值例子：$C = 10^{21}$时，用$N=10^9$训$D=C/(6N)$的loss vs 用$N=10^{10}$训$D=C/(6N)$的loss。

    \item \textbf{(更深的问题)} 第8讲说大batch训练存在"critical batch size" $B_{\text{crit}}$，超过它后加速比急剧下降。假设$B_{\text{crit}} \propto L / \|\nabla L\|^2$（梯度噪声尺度），对于loss更低的大模型，$B_{\text{crit}}$更大还是更小？这如何进一步支持"算力应给大模型"的论点？

    \item \textbf{(开放)} 以上分析假设所有数据质量相同。如果数据存在质量差异（如网页爬取 vs 精心标注），Scaling Law会如何修正？你认为"数据工程"在当前范式下的价值是什么？
\end{subproblems}
\end{problem}

%--------------------------------------------------
\subsection*{初始化：信号如何穿过深层网络？}
%--------------------------------------------------

\begin{problem}[$\star\star\star$ \textbf{从方差传播推导初始化方法}]
考虑$L$层全连接网络（无偏置，无BN）：$h^{(l)} = W^{(l)} \sigma(h^{(l-1)})$，其中$\sigma$是激活函数。每层宽度为$n$，$W^{(l)} \in \mathbb{R}^{n \times n}$，权重独立同分布于$\mathcal{N}(0, \sigma_w^2/n)$。

\begin{subproblems}
    \item \textbf{(前向方差传播)} 假设$h^{(0)}$的各分量独立、均值0、方差1。对于线性激活$\sigma(x) = x$，证明$\text{Var}[h_j^{(l)}] = (n \sigma_w^2)^l$（用归纳法）。要使信号不爆炸也不消失，$\sigma_w^2$应取多少？

    \item \textbf{(ReLU)} 对$\sigma(x) = \max(0, x)$，计算$\mathbb{E}[\sigma(z)^2]$（$z \sim \mathcal{N}(0,1)$）。由此证明前向方差传播变为$\text{Var}[h_j^{(l)}] = \left(\frac{n \sigma_w^2}{2}\right)^l$。要保持方差不变，$\sigma_w^2$应取多少？这就是\textbf{He初始化}。

    \item \textbf{(反向梯度传播)} 设loss对第$l$层输出的梯度为$\delta^{(l)}$。证明$\delta^{(l-1)} = (W^{(l)})^\top \delta^{(l)} \cdot \sigma'(h^{(l-1)})$。类似地分析$\text{Var}[\delta^{(l-1)}]$与$\text{Var}[\delta^{(l)}]$的关系。要使梯度不消失不爆炸，对$\sigma_w^2$的要求是什么？

    \item \textbf{(Xavier/Glorot的折中)} 前向要求$\sigma_w^2 = 1/n_{\text{in}}$，反向要求$\sigma_w^2 = 1/n_{\text{out}}$。当$n_{\text{in}} \neq n_{\text{out}}$时，Xavier初始化取$\sigma_w^2 = 2/(n_{\text{in}} + n_{\text{out}})$作为折中。解释这个选择的合理性。

    \item \textbf{(数值验证)} 取$L=50$, $n=256$。分别用$\sigma_w^2 = 1/n$（Xavier）、$\sigma_w^2 = 2/n$（He）、$\sigma_w^2 = 0.01/n$（太小）、$\sigma_w^2 = 10/n$（太大），计算第$L$层激活的理论方差。哪些会导致数值上溢或下溢？

    \item \textbf{(连接第6讲：正交初始化)} 如果$W^{(l)}$是正交矩阵（$W^\top W = I$），无论多少层，线性网络的前向和反向信号都不会衰减。解释这与Stiefel流形优化的联系。正交初始化在实践中与He初始化相比有什么优劣？

    \item \textbf{(开放)} Transformer中的初始化有何不同？残差连接$h^{(l)} = h^{(l-1)} + f^{(l)}(h^{(l-1)})$如何改变方差传播？如果有$L$个残差块，$f^{(l)}$的初始化尺度应该怎么随$L$调整？（提示：GPT-2使用$1/\sqrt{L}$缩放。）
\end{subproblems}
\end{problem}

%--------------------------------------------------
\subsection*{Batch Normalization的优化视角}
%--------------------------------------------------

\begin{problem}[$\star\star\star$ \textbf{BatchNorm到底在做什么？}]
Batch Normalization对每层的激活做归一化：$\hat{h}_j = (h_j - \mu_B)/\sigma_B$（$\mu_B, \sigma_B$是batch内的均值和标准差），再做仿射变换$\tilde{h}_j = \gamma \hat{h}_j + \beta$。

\begin{subproblems}
    \item \textbf{(第5讲的视角：预条件)} 考虑两层线性网络$y = w_2(w_1 x)$。如果放大$w_1$为$\alpha w_1$，同时缩小$w_2$为$w_2/\alpha$，输出不变。但$\nabla_{w_1} L$会变化。证明这种"尺度不变性破坏"使得loss landscape的条件数变差。BN通过归一化消除了激活的尺度——从这个角度解释BN为什么加速训练。

    \item \textbf{(Lipschitz平滑化)} 有研究表明BN的主要作用不是"减少internal covariate shift"，而是使loss landscape更加平滑（loss的梯度更Lipschitz）。回忆第5讲：$L$-光滑函数上GD的收敛率为$O(1/(LT))$。如果BN将$L$减小了一半，收敛率如何变化？这对步长选择有什么影响？

    \item \textbf{(BN与batch size的纠缠)} BN的统计量$\mu_B, \sigma_B$依赖于batch。当batch size很小时，估计不准；很大时统计量稳定。回忆第8讲的大batch训练，解释为什么BN在大batch（$\geq 4096$）下反而出现问题。Group Norm和Layer Norm是如何绕开这个问题的？

    \item \textbf{(开放)} 在当前的LLM训练中，几乎没有使用BatchNorm，而是用LayerNorm或RMSNorm。从"训练稳定性"和"并行友好度"两个角度分析原因。RMSNorm只保留了BN的哪一部分（均值中心化 vs 方差归一化）？为什么这就够了？
\end{subproblems}
\end{problem}

%--------------------------------------------------
\subsection*{Loss Landscape的几何与泛化}
%--------------------------------------------------

\begin{problem}[$\star\star\star$ \textbf{平坦极小值与泛化}]
\begin{subproblems}
    \item \textbf{(从第9讲出发)} 课程说高维loss landscape中大多数临界点是鞍点，低loss区域多是极小值。但极小值有"锐利"和"平坦"之分。直觉上解释为什么平坦的极小值（Hessian特征值小）可能泛化更好。

    \item \textbf{(PAC-Bayes视角)} PAC-Bayes界说：对参数加扰动$w \to w + \epsilon$（$\epsilon \sim \mathcal{N}(0, \sigma^2 I)$），如果训练loss不变，则泛化gap小。形式化地：
    \[
    L_{\text{test}}(w) \leq \mathbb{E}_\epsilon[L_{\text{train}}(w+\epsilon)] + O\left(\sqrt{\frac{\|w\|^2/(2\sigma^2) + \ln n}{n}}\right)
    \]
    解释为什么"平坦"区域对应更好的PAC-Bayes界。

    \item \textbf{(大batch vs 小batch)} 经验上，大batch SGD倾向于收敛到锐利极小值，小batch SGD倾向于收敛到平坦极小值。从SGD噪声的角度解释原因：SGD的有效噪声协方差为$\Sigma_{\text{noise}} \approx \frac{1}{B}\text{Cov}[\nabla f_i]$。$B$增大时噪声减小，为什么这会改变收敛到的极小值类型？

    \item \textbf{(SAM: Sharpness-Aware Minimization)} SAM的更新公式为：
    \[
    w \leftarrow w - \eta \nabla_w L(w + \rho \frac{\nabla L(w)}{\|\nabla L(w)\|})\bigg|_w
    \]
    即先在梯度方向"走一步"扰动，再计算该处的梯度做更新。解释SAM为什么倾向于找到平坦极小值。SAM每步的计算代价是普通SGD的多少倍？

    \item \textbf{(回到Scaling Law)} 如果大模型有更多的平坦极小值（过参数化使得loss landscape更"简单"），这为"大模型泛化更好"提供了另一个角度的解释。结合第9讲的过参数化隐式正则化理论，讨论模型规模与泛化的关系。

    \item \textbf{(开放)} "Double Descent"现象：随着模型大小增加，测试误差先降后升（经典偏差-方差），然后在过参数化区域再次下降。你认为这与loss landscape的几何有什么关系？过参数化区域的极小值有什么特殊性质？
\end{subproblems}
\end{problem}

%--------------------------------------------------
\subsection*{为什么深度有用？}
%--------------------------------------------------

\begin{problem}[$\star\star\star$ \textbf{深度的表达力与优化效率}]
\begin{subproblems}
    \item \textbf{(表达力)} 考虑用深度$L$、宽度$n$的ReLU网络逼近函数$f(x) = x^2$在$[0,1]$上。$L=1$（一层）时需要$n$个ReLU单元达到$O(1/n^2)$的逼近误差。证明$L=O(\log(1/\epsilon))$层、$n=O(1)$宽即可达到$\epsilon$误差。（提示：$x^2$可以通过逐步"对折"逼近。）

    \item \textbf{(深度矩阵分解)} 考虑矩阵补全问题：$\min_X \|P_\Omega(X - M)\|_F^2$ s.t. $\text{rank}(X) \leq k$。将$X$参数化为$X = W_L W_{L-1} \cdots W_1$（深度$L$的矩阵分解）。第9讲说深度分解有隐式正则化：
    \[
    \dot{\sigma}_r \propto \sigma_r^{2-2/L} \cdot \text{(gradient signal)}
    \]
    当$L$很大时，小奇异值$\sigma_r$的增长率远小于大奇异值。解释为什么这导致低秩偏好。

    \item \textbf{(连接第6讲)} 如果我们在Stiefel流形上优化$W_l$（正交约束），$\sigma_r$全部等于1，上述隐式正则化消失。这说明什么？为什么"不对称性"对深度学习的归纳偏置很重要？

    \item \textbf{(连接第9讲：NTK的局限)} NTK理论说无限宽网络等价于核方法——但核方法没有"深度优势"（特征不随训练变化）。这意味着有限宽度的特征学习（feature learning）才是深度的关键。从这个角度，NTK理论恰恰解释了"深度无用"的极限情形。讨论这对理解现实中深度网络的启示。

    \item \textbf{(开放)} Transformer的depth vs width trade-off：有研究表明增加深度比增加宽度更"计算高效"。从优化和表达力两个角度分析可能的原因。
\end{subproblems}
\end{problem}

%--------------------------------------------------
\subsection*{学习率与训练动力学}
%--------------------------------------------------

\begin{problem}[$\star\star\star$ \textbf{学习率的"相变"与Edge of Stability}]
\begin{subproblems}
    \item \textbf{(经典理论)} 第5讲说$L$-光滑函数上GD收敛要求步长$\eta < 2/L$（$L$是Hessian最大特征值）。证明：如果$\eta > 2/L$，则在Hessian最大特征值方向上迭代发散。

    \item \textbf{(Edge of Stability现象)} 但在实践中，全batch GD训练神经网络时观察到：Hessian最大特征值$\lambda_{\max}$会自动调整到$\approx 2/\eta$附近振荡！即系统"走到稳定性的边缘"但不崩溃。这违反了经典理论——经典理论假设$L$不变，但神经网络的$L$是会变的。

    从优化景观的角度解释：当$\lambda_{\max} > 2/\eta$时，最陡方向上的迭代会"弹开"，使参数移动到$\lambda_{\max}$较小的区域。这形成一种负反馈机制。画出$\lambda_{\max}$随训练的示意图。

    \item \textbf{(学习率Warmup)} 训练初期loss landscape曲率很大（$\lambda_{\max}$大），直接用大学习率会发散。Warmup在前$T_w$步线性增大学习率$\eta_t = \eta_{\max} \cdot t/T_w$。从Edge of Stability的视角解释warmup的作用：给网络时间让$\lambda_{\max}$先降下来。

    \item \textbf{(Cosine Decay)} 常见学习率衰减方案$\eta_t = \eta_{\min} + \frac{1}{2}(\eta_{\max} - \eta_{\min})(1 + \cos(\pi t / T))$。如果按第5讲的理论，凸函数上的最优步长衰减是$O(1/t)$或$O(1/\sqrt{t})$。Cosine Decay不属于任何一种，但在实践中更好。从"先探索后精化"的角度解释cosine schedule的设计直觉。

    \item \textbf{(连接第8讲)} 回忆SGD收敛率$O(\sigma^2\eta / (1-\eta L) + (1-\eta\mu)^T)$中的两项：优化误差和统计误差。学习率大时优化快但统计误差大。推导"两阶段"策略的最优切换点：先用大$\eta$快速降低优化误差，再用小$\eta$降低统计误差。这与warmup+decay的实践做法是否一致？

    \item \textbf{(开放)} 最近有研究尝试完全去掉学习率调度，用"Mechanistic Design"（如Muon优化器、Schedule-Free Adam）。你觉得学习率调度是"必须的物理"还是"工程补丁"？
\end{subproblems}
\end{problem}

%--------------------------------------------------
\subsection*{梯度裁剪与训练稳定性}
%--------------------------------------------------

\begin{problem}[$\star\star$ \textbf{梯度裁剪的理论基础}]
大模型训练中几乎必用梯度裁剪：$g \leftarrow g \cdot \min\left(1, \frac{c}{\|g\|}\right)$。
\begin{subproblems}
    \item \textbf{(第5讲视角)} 不裁剪时，梯度下降步$x \leftarrow x - \eta g$。如果$\|g\|$很大（远大于$L$-光滑性保证的范围），下降引理$f(x-\eta g) \leq f(x) - \eta\|g\|^2 + \frac{L\eta^2}{2}\|g\|^2$可能不再保证下降。解释梯度裁剪的作用：将有效步长限制为$\eta c / \|g\|$。

    \item \textbf{(连接第14讲：信任域)} 注意到裁剪后的更新满足$\|x_{k+1} - x_k\| \leq \eta c$——这正是一个信任域约束！从TRPO/PPO的思路理解梯度裁剪：它是最简单的信任域方法。

    \item \textbf{(实际考虑)} 在Transformer训练中，梯度裁剪阈值$c$通常设为$1.0$。如果$c$太小，训练太慢；太大，等于没裁剪。设计一个实验方案来找到最优的$c$。

    \item \textbf{(开放)} 除了梯度裁剪，还有哪些手段可以稳定大模型训练？（如loss scaling、gradient accumulation、架构层面的Pre-Norm vs Post-Norm。）从优化的角度将它们统一为"控制每步更新量"的不同策略。
\end{subproblems}
\end{problem}

%--------------------------------------------------
\subsection*{稀疏性与效率}
%--------------------------------------------------

\begin{problem}[$\star\star\star$ \textbf{Lottery Ticket与稀疏优化}]
\begin{subproblems}
    \item \textbf{(Lottery Ticket假说)} 一个随机初始化的稠密网络包含一个稀疏子网络（"中奖彩票"），该子网络\emph{使用原始初始化}可以训练到与稠密网络相同的精度。这意味着什么？从第9讲的过参数化角度：稠密网络的作用是否只是提供了足够多的"彩票"供优化选择？

    \item \textbf{(连接$\ell_1$正则化)} 第15讲的LASSO使用$\ell_1$范数促进稀疏。如果对网络权重加$\ell_1$正则化$\lambda\|w\|_1$，能找到"中奖彩票"吗？为什么简单的$\ell_1$在深度网络上效果不如"迭代剪枝"（train → prune → rewind → retrain）？

    \item \textbf{(Mixture of Experts)} 现代LLM（如Mixtral）使用稀疏MoE：每个token只激活$k/n$的参数。这是另一种"动态稀疏性"。从计算量的角度：如果MoE有$N$个参数但每个token只用$N/8$个，训练FLOPs接近$N/8$的稠密模型，但表达力接近$N$。结合Scaling Law讨论：MoE是否改变了$N$-$D$-$C$的最优分配？

    \item \textbf{(开放)} 人脑有$\sim 10^{11}$个神经元，但同一时刻只有$\sim 1\%$在活跃放电。讨论生物稀疏性与Lottery Ticket/MoE之间可能的联系。
\end{subproblems}
\end{problem}

%--------------------------------------------------
\subsection*{RLHF的优化本质}
%--------------------------------------------------

\begin{problem}[$\star\star\star$ \textbf{当奖励模型骗了你：Goodhart定律的优化解读}]
RLHF中，我们用PPO优化一个学习到的奖励模型$r_\phi(x,y)$（$x$是prompt，$y$是response）。
\begin{subproblems}
    \item \textbf{(第14讲出发)} PPO优化的目标是$\max_\theta \mathbb{E}_{x,y\sim\pi_\theta}[r_\phi(x,y)] - \beta D_{KL}(\pi_\theta \| \pi_{\text{ref}})$。如果没有KL惩罚（$\beta = 0$），会发生什么？$\pi_\theta$会如何利用$r_\phi$的漏洞？

    \item \textbf{(Reward Hacking)} 奖励模型$r_\phi$是在有限数据上训练的，在分布外（out-of-distribution）区域不可靠。当$\pi_\theta$偏离$\pi_{\text{ref}}$太远时，它可能找到$r_\phi$很高但人类评价很低的response。画出示意图：横轴为$D_{KL}(\pi_\theta \| \pi_{\text{ref}})$，纵轴分别为$r_\phi$（proxy reward）和真实人类偏好（true reward）。两条曲线在什么区域开始分叉？

    \item \textbf{(连接第2讲：约束优化)} 从KKT的角度重新理解$\beta$：它是KL约束$D_{KL} \leq \delta$的对偶变量。$\beta$太小$\Leftrightarrow$约束太松$\Leftrightarrow$reward hacking。$\beta$太大$\Leftrightarrow$约束太紧$\Leftrightarrow$策略几乎不更新。最优的$\beta$对应什么？（互补松弛条件意味着什么？）

    \item \textbf{(DPO: 绕开RL)} Direct Preference Optimization直接优化偏好数据，不需要训练显式的奖励模型，也不需要PPO。其关键观察是：最优策略$\pi^*$和奖励$r$之间有闭式关系$r(x,y) = \beta \log \frac{\pi^*(y|x)}{\pi_{\text{ref}}(y|x)} + \beta \log Z(x)$。因此可以将$r$消去，直接用偏好对$(y_w \succ y_l)$优化$\pi$。

    从优化复杂度的角度比较PPO和DPO：PPO需要4个模型（policy, reference, reward, value），DPO只需要2个（policy, reference）。这对GPU显存意味着什么？

    \item \textbf{(开放)} 你认为"对齐税"（alignment tax）——即对齐后模型在基准测试上的能力下降——是不可避免的吗？从多目标优化的帕累托前沿角度分析"能力"和"安全"之间的trade-off。
\end{subproblems}
\end{problem}

%--------------------------------------------------
\subsection*{注意力机制的优化与效率}
%--------------------------------------------------

\begin{problem}[$\star\star\star$ \textbf{Attention为什么有效，以及如何加速它}]
自注意力机制：$\text{Attn}(Q,K,V) = \text{softmax}(QK^\top / \sqrt{d_k})V$，其中$Q,K,V \in \mathbb{R}^{n \times d_k}$。
\begin{subproblems}
    \item \textbf{(第7讲视角：低秩近似)} 注意力矩阵$A = \text{softmax}(QK^\top/\sqrt{d_k})$是$n \times n$的。如果$A$近似低秩（少数几个token主导注意力），那么$AV$可以用随机化方法近似。设$A$的有效秩为$k \ll n$，估计近似计算的加速比。

    \item \textbf{(IO复杂度)} 第16讲的FlashAttention的核心洞察不是减少FLOPs，而是减少HBM（高带宽内存）访问。标准Attention需要将$n \times n$的注意力矩阵写入HBM再读出做$\times V$。FlashAttention通过分块计算（tiling），在SRAM中完成softmax和$\times V$，只写最终结果到HBM。

    假设SRAM大小$M$，矩阵$Q,K,V$各$n \times d$。标准Attention的HBM访问量为$O(nd + n^2)$，FlashAttention为$O(n^2 d^2/M)$。当$d^2 < M$（通常成立，$d=64$, $M \sim 10^5$）时，FlashAttention的IO复杂度近似$O(n^2 d^2/M)$。与标准的$O(n^2)$比较，加速因子是什么？

    \item \textbf{(连接第4讲：分治)} KV-Cache在自回归生成中缓存已计算的$K,V$。生成第$t$个token时，只需计算新的$q_t$与所有历史$K_{1:t}$的注意力。这将生成$n$个token的复杂度从$O(n^3 d)$降到$O(n^2 d)$。推导这个结果。

    \item \textbf{(稀疏注意力)} 如果限制每个token只关注$k$个最近的token（sliding window）加上$s$个全局token，注意力矩阵变成稀疏的。计算FLOPs从$O(n^2 d)$降到多少？这对长文本（$n > 100000$）意味着什么？

    \item \textbf{(开放)} 近年来有很多"线性注意力"工作（如Mamba, RWKV），将$O(n^2)$降到$O(n)$。但它们在某些任务上不如标准Transformer。你认为$O(n^2)$的注意力捕获了什么是$O(n)$架构难以捕获的？从"表达力 vs 效率"的trade-off讨论。
\end{subproblems}
\end{problem}

%--------------------------------------------------
\subsection*{涌现、相变与优化}
%--------------------------------------------------

\begin{problem}[$\star\star\star$ \textbf{为什么大模型会"涌现"新能力？}]
\begin{subproblems}
    \item \textbf{(相变的类比)} 第1讲讨论了计算复杂性中的相变（如随机SAT在子句密度$m/n \approx 4.27$处从几乎全可满足变为几乎全不可满足）。类似地，LLM在模型规模增大时，某些能力（如思维链推理、多位数加法）似乎"突然出现"。

    从优化的角度思考：如果某个能力需要网络学到一个特定的"电路"（子网络），而这个电路需要最少$N_{\min}$个参数才能表达，那么在$N < N_{\min}$时loss在该任务上无法下降，$N > N_{\min}$时loss突然可以下降。画出"某任务的loss vs 模型大小"的示意图。

    \item \textbf{(Grokking)} 另一个"相变"现象：网络先记住训练数据（过拟合），很久之后突然学会泛化规律（grokking）。从隐式正则化的角度：GD的轨迹先到达一个"记忆解"（高复杂度），然后缓慢漂移到"泛化解"（低复杂度）。这与第9讲的最小范数偏好有什么联系？

    \item \textbf{(连接第11讲：探索-利用)} 如果将训练过程类比为Bandit问题：模型在"简单特征"和"复杂特征"之间分配"学习容量"。训练初期利用简单特征就能降loss（exploitation），后期需要发现更复杂的模式（exploration）。学习率调度是否可以理解为控制这种探索-利用的balance？

    \item \textbf{(开放)} "涌现"是真实的还是度量（metric）的假象？有研究指出如果用连续的loss（而非accuracy这样的阈值度量），"涌现"变成了平滑的改进。你如何设计实验来区分"真正的相变"和"度量假象"？
\end{subproblems}
\end{problem}

%--------------------------------------------------
\subsection*{Inference优化}
%--------------------------------------------------

\begin{problem}[$\star\star$ \textbf{推理比训练更重要：推理优化的数学}]
训练一个LLM花费$C_{\text{train}}$，但部署后每个query的推理成本$c_{\text{infer}}$会被调用$Q$次。总成本$C_{\text{total}} = C_{\text{train}} + Q \cdot c_{\text{infer}}$。
\begin{subproblems}
    \item 当$Q > C_{\text{train}} / c_{\text{infer}}$时，推理成本主导。对GPT-4量级（$C_{\text{train}} \sim 10^{25}$ FLOPs, $Q \sim 10^{12}$ queries），估算$c_{\text{infer}}$的临界值。

    \item \textbf{(量化)} 将FP16权重量化到INT4，理论上推理速度提升多少？（考虑带宽瓶颈：自回归生成是memory-bandwidth bound，速度$\propto$模型大小/内存带宽。）权重从2字节变为0.5字节，速度提升约$4\times$。但量化带来精度损失。从第7讲的随机化视角：量化可以看作一种"有偏sketching"。如何用随机舍入（stochastic rounding）减小偏差？

    \item \textbf{(Speculative Decoding)} 用一个小模型（draft model）快速生成$k$个候选token，再用大模型（target model）并行验证。如果小模型的接受率为$\alpha$，期望每次验证通过的token数为$k\alpha$。总加速比近似$k\alpha / (c_{\text{draft}} \cdot k + c_{\text{target}})$。如果$c_{\text{draft}} = c_{\text{target}}/10$, $\alpha = 0.8$, $k=5$，计算加速比。

    \item \textbf{(连接第4讲)} 自回归生成本质上是顺序的（第$t$个token依赖第$t-1$个）。这与第4讲中"贪心算法的GPU困境"类似。Speculative Decoding是如何"打破"这种顺序依赖的？从分治法的"投机执行"（speculative execution）角度理解。

    \item \textbf{(开放)} "Test-time compute"（如OpenAI o1的思维链推理）用更多推理计算换取更好的回答。从在线优化的角度：这相当于在推理时做额外的"搜索"。你认为训练计算和推理计算之间存在Scaling Law吗？
\end{subproblems}
\end{problem}

\end{document}
